\documentclass[11pt,a4paper]{article}

% =============================================================================
% additional_results.tex -- companion volume to preliminary_results.tex
% =============================================================================
% Extends and complements the headline memo (preliminary_results.tex) with
% deeper descriptive cuts, per-firm decompositions, and the regulatory /
% market-structure context that the headline memo only summarises. Where
% this document analyses the same outcomes as the preliminary memo it
% follows the same methodology (Spec A: BSTS for prices and cleared MW;
% Spec B: per-hour-class BSTS on in-band offered energy; Spec C: per-curve
% DiD for bid shape) and cross-references rather than re-derives.
%
% Renamed from descriptive_facts.tex 2026-05-31.
% Content originally migrated 2026-05-17 from thesis/paper/provisional.tex.
% The regression-style identification content lives in
% thesis/provisional/regression_results.tex.

\PassOptionsToPackage{table}{xcolor}
\usepackage{amsmath, amssymb, amsthm, mathtools}
\usepackage{dsfont}
\usepackage{graphicx, booktabs, subcaption, tikz, float}
\graphicspath{{../../figures/thesis/}}
\usepackage{array, makecell, multirow}
\usepackage{xcolor}
\usepackage{longtable}
\usepackage[hidelinks]{hyperref}
\usepackage{geometry}
\geometry{a4paper, margin=2.5cm}
\usepackage{pdflscape}
\usepackage{caption}
\captionsetup{font=small, labelfont=bf}

% Color for seasonality-adjusted values shown alongside raw values in tables.
\definecolor{seasoncol}{rgb}{0.10, 0.35, 0.65}

% \code{...} -- monospace path with line-break-at-/-_-. (no \_ escaping needed)
\usepackage{url}
\urlstyle{tt}
\Urlmuskip=0mu plus 1mu\relax
\def\UrlBreaks{\do\/\do\_\do\-\do\.\do\:\do\,}
\newcommand{\code}[1]{\nolinkurl{#1}}

\title{Additional results\\[3pt]\large Companion volume to \emph{preliminary results}: deeper descriptive cuts, per-firm decompositions, and market-structure context}
\author{Pablo Paramio P\'erez}
\date{\today}

\begin{document}
\maketitle

\noindent\emph{Extended robustness and other findings to \textbf{preliminary results}: Spec~C extensions (Cogen / Biomass / Nuclear in \S\ref{sec:prov:cogen_stability}; per-firm $\times$ tech in \S\ref{sec:prov:per_firm_tech_shape}), plus findings outside the headline memo's scope (price-setting, geographic CCGT dominance, post-DA intervention chain, cost cascade). All identification claims use the same methodology (Spec~A BSTS, Spec~B per-hour-class BSTS, Spec~C per-curve DiD), same window-and-market-specific bandwidth $h$ ($p_{90}-p_{50}$ of MCP per cell; \emph{preliminary results} \S 1), same 2024 placebo. Seasonality handled by BSTS internally plus same-calendar matching; per-curve metrics at the wider $h = 140$ are descriptive supplements only.}


\tableofcontents

\bigskip

% =====================================================================
\section*{Glossary of acronyms}\label{sec:bid_internal:glossary}
% =====================================================================
\addcontentsline{toc}{section}{Glossary of acronyms}

\noindent\footnotesize
\textbf{CCGT} combined-cycle gas turbine ---
\textbf{DA} day-ahead market ---
\textbf{IDA} intraday auctions (3 European SIDC sessions post-2024-06-14) ---
\textbf{MCP} market clearing price ---
\textbf{MTU} market time unit (MTU60 hourly, MTU15 quarter-hourly) ---
\textbf{ISP15} Imbalance Settlement Period at 15 min (started 2024-12-01) ---
\textbf{REE} Red Eléctrica de España (TSO) ---
\textbf{OMIE} wholesale market operator ---
\textbf{PDBC} OMIE DA-cleared schedule ---
\textbf{PDBF} PDBC plus bilaterals ---
\textbf{PHF} post-IDA + post-RT final schedule ---
\textbf{TR} Tiempo Real real-time restrictions (PO 3.4--3.5, pay-as-bid) ---
\textbf{Fase I/II} pre-IDA / post-IDA technical restrictions (pay-as-bid) ---
\textbf{aFRR / mFRR / RR} balancing reserves (PICASSO / MARI / TERRE, marginal) ---
\textbf{BSP} Balance Service Provider ---
\textbf{EUPHEMIA} DA pan-European auction algorithm.\normalsize

\bigskip

% =====================================================================
\section*{Reform-window calendar and confounders}\label{sec:bid_internal:timeline}
% =====================================================================
\addcontentsline{toc}{section}{Reform-window calendar and confounders}

\noindent We use the same five-regime reform-window timeline as the headline memo (\emph{preliminary results} \S 1, Figure 1). Two confounders to flag for any cross-regime reading in this document.

\textbf{Calendar non-overlap.} The five regime windows cover different calendar months (3-sess Jun--Nov 2024; ISP15-win Dec 2024--Mar 2025; MTU15-IDA pre-blk late Mar--Apr 2025; MTU15-IDA post-blk May--Sep 2025; DA15/ID15 Oct 2025 onwards). Spanish electricity is strongly seasonal, so raw cross-regime means mix the regime effect with the seasonal mix. Three regimes (3-sess, ISP15-win, MTU15-IDA pre-blk) have a clean same-calendar pre-IDA baseline available (Jun--Nov, Dec--Mar, late-Mar--Apr in 2018--2023, none under any granularity reform); the post-blk and DA15/ID15 windows are reforzada-confounded with no clean granularity-only same-cal comparator. \emph{Identification claims defer to the headline memo's BSTS + 2024 same-calendar placebo, where seasonality is absorbed by construction.}

\textbf{Cross-border co-treatment for ID15.} MTU15-DA was an EU-wide ``Big Bang'' on 2025-09-30: all SDAC bidding zones cleared their first 15-min day-ahead simultaneously, so cross-border DA allocation changed identically for Spain and its neighbours --- a critical/flat DiD on cross-border net flow at the cutover is a clean null ($\theta = 793$ MW, $t = 0.73$). MTU15-IDA is different: it cut the bid time unit of the already-pan-European SIDC and XBID venues from 60 to 15 min, so cross-border \emph{allocation timing} co-treats with bid-granularity --- the same DiD on net flow returns $\theta = 8\,965$ MW ($t = 5.12$), persistent across post-window truncations. The DA15 cross-border channel is therefore empirically inert; any ID15 headline reads as bundling bid-granularity with cross-border-allocation-timing within the same intervention.\footnote{SDAC 15-minute MTU Information Package (\url{https://www.nemo-committee.eu/assets/files/sdac-15-minute-mtu-information-paper.pdf}). Cross-border flows: ESIOS indicators 535 (FR import), 536 (FR export), 539 (PT import), 540 (PT export); net flow per (date, clock-hour) computed as $(535+539)-(536+540)$ summed across the 4 quarters of the hour.}

\bigskip

% =====================================================================
\section*{BSTS / CausalImpact: model, math and configurations}\label{sec:bid_internal:bsts_methods}
% =====================================================================
\addcontentsline{toc}{section}{BSTS / CausalImpact: model, math and configurations}

\noindent Every Spec~A result in this document --- and every BSTS effect cited from \emph{preliminary results} \S 4.A --- is produced by \texttt{CausalImpact::CausalImpact} (Brodersen \emph{et al.} 2015), an R wrapper around the Bayesian structural time-series engine \texttt{bsts} of Scott \& Varian (2014). The model is a linear-Gaussian state-space form fit by Markov-chain Monte Carlo on the pre-intervention window and projected as a posterior counterfactual into the post-intervention window.

\subsection*{Notation and structural form}

Let $y_t \in \mathbb{R}$ be the outcome at time $t = 1, \dots, T$, and $x_t \in \mathbb{R}^p$ a vector of contemporaneously-observed covariates. The intervention time $T_0$ partitions the index into a pre-window $\mathcal{T}_0 = \{1,\dots,T_0\}$ and a post-window $\mathcal{T}_1 = \{T_0{+}1,\dots,T\}$. The data-generating process for $y_t$ decomposes into a trend, $K$ seasonal cycles, a contemporaneous regression, and Gaussian observation noise:
\begin{equation}
y_t \;=\; \mu_t \;+\; \sum_{k=1}^{K} S_t^{(k)} \;+\; \beta^{\!\top} x_t \;+\; \varepsilon_t,
\qquad \varepsilon_t \stackrel{\text{iid}}{\sim} \mathcal{N}\!\bigl(0,\, \sigma_y^2\bigr).
\label{eq:bsts_dgp}
\end{equation}

\subsection*{Components}

\noindent\textbf{Local level (trend).} The trend is a random walk \emph{without} drift slope:
\begin{equation}
\mu_{t+1} \;=\; \mu_t \;+\; \eta_t^{\mu},
\qquad \eta_t^{\mu} \stackrel{\text{iid}}{\sim} \mathcal{N}\!\bigl(0,\, \sigma_\mu^2\bigr).
\label{eq:local_level}
\end{equation}
The local linear trend variant (slope state $\nu_t$, $\mu_{t+1} = \mu_t + \nu_t + \eta_t^\mu$, $\nu_{t+1} = \nu_t + \eta_t^\nu$) is rejected: a slope state projects linearly into $\mathcal{T}_1$, extrapolating years of pre-window growth into the post-window. The local level lets the underlying $y$-level wander without imposing a direction.

\medskip
\noindent\textbf{Dummy seasonal of integer period $S$.} Constrained-sum recursion:
\begin{equation}
S_t^{(\text{dum},S)} \;=\; -\sum_{j=1}^{S-1} S_{t-j}^{(\text{dum},S)} \;+\; \omega_t^{(S)},
\qquad \omega_t^{(S)} \stackrel{\text{iid}}{\sim} \mathcal{N}\!\bigl(0,\, \sigma_S^2\bigr).
\label{eq:dummy_seasonal}
\end{equation}
The state vector holds the last $S{-}1$ values of $S_t^{(\text{dum},S)}$; the constraint $\sum_{j=0}^{S-1} S_{t-j}^{(\text{dum},S)}$ is small (only $\omega_t$-wide), so the seasonal averages to zero over each full cycle. Standard choices: weekly ($S=7$) and diurnal ($S=24$).

\medskip
\noindent\textbf{Trigonometric (Harvey) seasonal of period $P$ and $F$ harmonics.} For each harmonic $f \in \{1,\dots,F\}$ with frequency $\lambda_f = 2\pi f / P$, define a pair of states $(\gamma_{f,t}, \gamma_{f,t}^{\ast}) \in \mathbb{R}^2$ that rotates with innovation:
\begin{equation}
\begin{pmatrix} \gamma_{f,t+1} \\[2pt] \gamma_{f,t+1}^{\ast} \end{pmatrix}
\;=\;
\begin{pmatrix} \cos\lambda_f & \sin\lambda_f \\[2pt] -\sin\lambda_f & \cos\lambda_f \end{pmatrix}
\begin{pmatrix} \gamma_{f,t} \\[2pt] \gamma_{f,t}^{\ast} \end{pmatrix}
\;+\;
\begin{pmatrix} \eta_{f,t} \\[2pt] \eta_{f,t}^{\ast} \end{pmatrix},
\qquad \eta_{f,t},\, \eta_{f,t}^{\ast} \stackrel{\text{iid}}{\sim} \mathcal{N}\!\bigl(0,\, \sigma_{\text{trig}}^2\bigr).
\label{eq:trig_seasonal}
\end{equation}
The contribution to $y_t$ is the sum of the ``real'' components: $S_t^{(\text{trig},P,F)} = \sum_{f=1}^F \gamma_{f,t}$. The block costs $2F$ states regardless of $P$. For an annual cycle on daily data ($P = 365.25$) with $F = 2$, this is $4$ states, against $364$ for the dummy form.

\medskip
\noindent\textbf{Regression with spike-and-slab.} The covariate vector enters linearly, with each coefficient carrying a binary inclusion indicator $\rho_j \in \{0,1\}$ and a Gaussian slab:
\begin{equation}
\beta_j \;=\; \rho_j\, \tilde\beta_j,
\qquad
\rho_j \stackrel{\text{iid}}{\sim} \mathrm{Bernoulli}(\pi),
\qquad
\tilde\beta \mid \rho \;\sim\; \mathcal{N}_{|\rho|}\!\bigl(b_\rho,\, V_\rho\bigr),
\label{eq:spike_slab}
\end{equation}
with $|\rho| = \sum_j \rho_j$. The slab covariance $V_\rho$ is a Zellner-style $g$-prior calibrated to $X^{\!\top}X$; the prior inclusion probability is $\pi = E[m]/p$ with default expected model size $E[m] = 3$.

\subsection*{State-space canonical form}

Collecting the dynamic components into a state vector $\alpha_t \in \mathbb{R}^d$ (local level + dummy seasonals + trig pairs, with $\beta$ held separately as a static parameter), the model is the linear-Gaussian state-space system
\begin{align}
\alpha_{t+1} &\;=\; T\, \alpha_t + R\, \eta_t,
& \eta_t &\stackrel{\text{iid}}{\sim} \mathcal{N}(0,\, Q),
\label{eq:ssm_state}\\
y_t &\;=\; Z\, \alpha_t + \beta^{\!\top} x_t + \varepsilon_t,
& \varepsilon_t &\stackrel{\text{iid}}{\sim} \mathcal{N}(0,\, \sigma_y^2),
\label{eq:ssm_obs}
\end{align}
where the transition matrix $T$ is block-diagonal,
\[
T \;=\; \mathrm{blockdiag}\!\Bigl(\,
\underbrace{1}_{\text{level}},\;
\underbrace{T_{\text{dum},S}}_{(S-1)\times(S-1)\ \text{companion}},\;
\underbrace{T_{\text{trig},F,P}}_{2F\times 2F\ \text{block-rotation in (\ref{eq:trig_seasonal})}}
\,\Bigr),
\]
$R$ is the selector matrix that picks out the components receiving innovations, $Z$ is the time-invariant observation row that extracts $\mu_t$, the leading dummy-seasonal slot, and $\sum_f \gamma_{f,t}$, and the innovation covariance is $Q = \mathrm{diag}\!\bigl(\sigma_\mu^2,\, \sigma_S^2 \mathbf{1}_{S-1},\, \sigma_{\text{trig}}^2 \mathbf{1}_{2F}\bigr)$.

\subsection*{Priors}

\noindent All innovation variances and the observation variance carry conjugate inverse-Gamma priors:
\[
\sigma_\mu^2,\, \sigma_S^2,\, \sigma_{\text{trig}}^2,\, \sigma_y^2
\;\stackrel{\text{indep}}{\sim}\; \mathrm{InvGamma}\!\bigl(\nu_{\bullet}/2,\, s_{\bullet}/2\bigr),
\]
with default prior degrees of freedom $\nu_{\bullet} \approx 1$ and prior scales $s_{\bullet}$ calibrated to the empirical variance of $y$. The regression prior is (\ref{eq:spike_slab}). All priors are mildly informative; at the daily/hourly precision used here, the posterior is dominated by the likelihood.

\subsection*{Posterior inference (Gibbs sampling)}

\noindent The joint posterior $p\!\bigl(\alpha_{1:T_0},\, \theta \,\big|\, y_{1:T_0}\bigr)$, with $\theta = (\sigma_\mu^2,\, \sigma_S^2,\, \sigma_{\text{trig}}^2,\, \sigma_y^2,\, \beta,\, \rho)$, has no closed form. We sample from it via a two-block Gibbs scheme:
\begin{enumerate}\setlength{\itemsep}{2pt}\setlength{\topsep}{2pt}
\item[\textbf{(a)}] \emph{States given parameters.} For fixed $\theta$, draw $\alpha_{1:T_0} \sim p(\cdot \mid y_{1:T_0},\, \theta)$ via the Durbin--Koopman simulation smoother: forward Kalman filter, backward simulation pass.
\item[\textbf{(b)}] \emph{Parameters given states.} Conditional on $\alpha_{1:T_0}$, the innovation variances draw from inverse-Gamma conjugate posteriors; $(\beta,\rho)$ draws from a stochastic-search variable-selection (SSVS) update that integrates out the slab and toggles each $\rho_j$ by a Bernoulli step.
\end{enumerate}
We run \texttt{niter} $= 2000$ iterations with the default $10\%$ burn-in; convergence is monitored via the effective sample size reported by \texttt{bsts}.

\subsection*{Counterfactual and treatment-effect summaries}

\noindent The post-window counterfactual is the posterior predictive of $y_t$ under the trained model evaluated at the \emph{actually-observed} post-window covariates $x_t$ (no extrapolation):
\[
\hat y_t^{(0)} \;\sim\; p\!\bigl(y_t \mid y_{1:T_0},\, x_t\bigr), \qquad t \in \mathcal{T}_1.
\]
The pointwise treatment effect is $\delta_t = y_t - \hat y_t^{(0)}$. We report the average effect
\[
\bar\delta \;=\; \frac{1}{|\mathcal{T}_1|}\!\sum_{t \in \mathcal{T}_1}\! \delta_t,
\]
its $95\%$ credible interval, and the Bayesian tail $\Pr\!\bigl(\bar\delta \le 0 \,\big|\, y_{1:T}\bigr)$ (the BSTS ``p-value'', a posterior tail, not a frequentist test).

\paragraph{What is extrapolated, and the robustness check for what isn't.} In the post-window the local level $\mu_t$ is held flat at $\mu_{T_0}$; the dummy seasonal cycles continue forward at their pre-window-fitted shapes; the regression contribution $\beta^{\!\top} x_t$ is evaluated at the \emph{actually-observed} post-window covariates. Annual seasonality therefore enters the counterfactual primarily through the covariate block: wind, solar and gas all carry strong annual cycles and their post-window realisations are observed --- so cross-month level shifts (winter heating, summer cooling, solar amplitude) are carried by the regression term, not by extrapolating any seasonal state. Residual annual variation that (wind, solar, gas) do not span is the one channel exposed to bias when the pre and post windows fall on different calendar bands. The placebo BSTS at fake cutover 2024-10-01 (\emph{preliminary results} \S 4.A.6) is the robustness check for this residual leakage --- a non-zero placebo $\bar\delta$ would wound the corresponding real estimate. The multi-cycle extension above is held in reserve as the heavier remedy if the placebo were to flag something.

\subsection*{Cycle stack: baseline and proposed extension}

\noindent\textbf{Baseline (current implementation).} Local level + one dummy seasonal + spike-and-slab regression on (wind, solar, gas). The seasonal block is weekly ($S = 7$, $6$ states) on daily data, or diurnal ($S = 24$, $23$ states) on hourly data. The state dimension is small ($d \in [7, 24]$) and the one-line \texttt{model.args} call to \texttt{CausalImpact} accommodates it. Low-frequency cross-month variation that the weekly/diurnal cycle cannot absorb is left to the regression block on (wind, solar, gas).

\medskip
\noindent\textbf{Proposed multi-cycle extension (not yet run).} Three seasonal layers plus the local level, to disentangle short-, medium- and low-frequency seasonality explicitly:
\begin{center}
\renewcommand{\arraystretch}{1.2}
\setlength{\tabcolsep}{8pt}
\footnotesize
\begin{tabular}{l l l c}
\toprule
\textbf{Layer} & \textbf{Form} & \textbf{Configuration} & \textbf{States} \\
\midrule
Diurnal (hourly data only) & dummy & $S = 24$, one obs per slot                & $23$ \\
Weekly                     & dummy & $S = 7$, one day (or $24$ hours) per slot & $6$  \\
Annual                     & trig. & $P = 365.25$ d (or $8766$ hr), $F = 2$ harmonics & $4$ \\
\bottomrule
\end{tabular}
\end{center}
The annual trig layer absorbs the systematic Apr--Sep vs Oct--Mar level differences in Spanish electricity (winter heating, summer cooling, hydro inflows, solar amplitude) that the baseline weekly cycle cannot reach and that the existing implementation leaves to the regression block. Adopting it requires constructing the \texttt{bsts} state specification manually --- \texttt{AddLocalLevel} $\to$ \texttt{AddSeasonal} $\to$ \texttt{AddTrig} --- and passing the resulting model to \texttt{CausalImpact} via the \texttt{bsts.model} interface; the existing one-line \texttt{model.args} call exposes only a single seasonal layer.

\medskip
\noindent\textbf{Why trigonometric, not dummy, for the annual layer.} A $365$-slot dummy seasonal needs three full years of pre-window before each slot has even three independent draws; identification is poor and MCMC wastes iterations on slots the data barely informs. The trigonometric form imposes smooth low-rank structure at $2F$ states regardless of period $P$. The harmonic count $F$ controls bias--variance: $F = 1$ gives a single sinusoid; $F = 2$ admits a winter \emph{and} a summer peak with an autumn dip --- the natural shape for Spanish electricity, where December heating and July--August cooling are both demand peaks.

\subsection*{Pre/post windows}

\noindent\textbf{MTU15-IDA:} pre 2022-01-01 to 2025-03-18, post 2025-03-19 to 2025-04-27 ($40$ days, pre-blackout, reforzada-absent). \textbf{MTU15-DA $\alpha$:} pre 2025-04-28 to 2025-09-30 ($5$ months, post-reforzada-onset, pre-MTU15-DA), post 2025-10-01 to 2025-12-31 --- separates granularity from reforzada. \textbf{MTU15-DA $\beta$:} pre 2022-01-01 to 2025-09-30 (long pre, joint granularity + reforzada). The spatial blackout-cutover BSTS in \S\ref{sec:prov:geo:spatial_blackout} uses pre 2024-06-14 to 2025-04-27, post 2025-04-28 to 2026-01-31.

\subsection*{Why BSTS over a panel DiD for level outcomes}

\noindent The DA price and per-tech cleared MWh series are strongly non-stationary (Spanish renewable capacity grew $\sim 6\times$ over 2018--2025) and carry a multi-scale seasonal structure. A panel DiD on a single outcome series cannot absorb all the cycles plus the regression transparently. BSTS decomposes the variation explicitly via (\ref{eq:bsts_dgp}), lets the spike-and-slab select among candidate covariates, and produces a counterfactual conditioned on the actually-realised post-period exogenous variables --- not on out-of-sample extrapolation. The DiD framework is reserved here for the within-day redistribution objects (Spec~C $\sigma_p$, $N_{\text{eff}}$), where critical and flat curves clear within the same day and the comparison absorbs all between-day variation by construction --- a different identification logic where BSTS would lose precision.

\bigskip

\noindent\textbf{Reading guide.} Part A is about bidding and trading patterns. Part B is about the regulatory/post-clearing channels (reforzada, geographic dominance, CNMC sanctions) that are independent of the granularity reform but constrain how the identification can be set up.

% =============================================================================
\part{Part A --- Bidding and trading patterns}
% =============================================================================

\noindent\emph{Reading guide for Part A.} Three setup sections come first. \S\ref{sec:prov:hour_class} defines the \textbf{hour-class taxonomy} (critical / flat / midday) used to slice every table and figure that follows. \S\ref{sec:prov:pricesetter} characterises \textbf{who price-sets} in the day-ahead market (\S\ref{sec:prov:pricesetter_da}) and, methodologically, in the intraday markets (\S\ref{sec:prov:pricesetter_ida}). \S\ref{sec:prov:dw_decomp} reports the bid-shape metrics ($\sigma_p$, $N_{\text{eff}}$ --- the same per-curve functionals as the headline memo \emph{preliminary results} \S 1, \S 4.C). Quantitative tables come in \S\ref{sec:prov:dw_tables}; the Spec~C extensions at the headline-memo bandwidth follow in \S\ref{sec:prov:per_firm_tech_shape} (per-firm $\times$ tech) and \S\ref{sec:prov:cogen_stability} (tech-scope robustness).

% =====================================================================
\section{Hour-class taxonomy --- supplementary classifier table}\label{sec:prov:hour_class}
% =====================================================================

\noindent The critical/flat partition is defined and motivated in
\emph{preliminary results} \S 2 (which uses the morning-ramp $\cup$
evening-ramp union as ``critical'', overnight hours as ``flat'',
midday as a separate falsification class). This subsection only
supplies the full per-hour classifier table that is too long for the
headline memo, plus two validation checks.

\begin{table}[H]
\centering\footnotesize
\caption{\textbf{Hour-of-day classification metrics, Spain 2023--2026} (ENTSO-E A65 + A75, 15-min ISP, $n \approx 3.4$ years of daily observations per hour). The classifier is $\sigma_{\text{within}}$, the expected SD of residual demand (load $-$ wind $-$ solar) across the 4 quarters of each clock-hour. Critical hours: $\sigma_{\text{within}} \in [522, 1303]$. Flat hours: $\sigma_{\text{within}} \in [336, 357]$. Midday: $\sigma_{\text{within}} \in [468, 502]$ (intermediate; excluded from main specs because of the solar-saturation channel). Dropped (transition) hours sit at the regime boundaries. \textit{Year-stability:} across 2023, 2024, 2025 the critical-set min $\sigma_{\text{within}}$ exceeds the flat-set max by $1.4\times$; no within-class swaps across years.}
\label{tab:hour_classification_full}
\newcommand{\critrow}{\rowcolor{red!10}}%
\newcommand{\flatrow}{\rowcolor{blue!10}}%
\begin{tabular}{r r r r r r l}
\toprule
$h$ & load (GW) & solar (GW) & net-load (GW) & $\sigma_{\text{within}}$ (MW) & $\Delta$load (MW) & class \\
\midrule
{ 0} & 23.2 & 0.2 & 16.1 & 398 & -864 & \textit{dropped} \\
\textit{ 1} & 22.4 & 0.1 & 15.4 & 356 & -481 & \textit{flat} \\
\textit{ 2} & 22.0 & 0.1 & 15.1 & 272 & -193 & \textit{flat} \\
\textit{ 3} & 21.9 & 0.1 & 15.2 & 328 & +194 & \textit{flat} \\
{ 4} & 22.9 & 0.1 & 16.3 & 654 & +1179 & \textit{dropped} \\
\textbf{ 5} & 24.9 & 0.3 & 18.4 & 768 & +1795 & \textit{critical} \\
\textbf{ 6} & 27.0 & 1.9 & 18.9 & 1004 & +1420 & \textit{critical} \\
\textbf{ 7} & 28.3 & 5.8 & 16.7 & 1261 & +602 & \textit{critical} \\
\textbf{ 8} & 28.7 & 10.5 & 12.9 & 1122 & +218 & \textit{critical} \\
{ 9} & 28.8 & 13.5 & 10.2 & 795 & -102 & \textit{dropped} \\
{10} & 28.6 & 15.0 & 8.7 & 583 & -78 & \textit{dropped} \\
{11} & 28.7 & 15.5 & 8.2 & 481 & +116 & \textit{dropped} \\
{12} & 28.6 & 15.5 & 8.0 & 438 & -60 & \textit{dropped} \\
{13} & 28.3 & 15.2 & 7.9 & 458 & -290 & \textit{dropped} \\
{14} & 28.0 & 14.4 & 8.2 & 550 & -57 & \textit{dropped} \\
{15} & 28.3 & 13.0 & 9.4 & 758 & +344 & \textit{dropped} \\
\textbf{16} & 28.9 & 10.1 & 12.5 & 1237 & +634 & \textit{critical} \\
\textbf{17} & 30.0 & 6.4 & 16.7 & 1367 & +902 & \textit{critical} \\
\textbf{18} & 31.1 & 2.6 & 21.3 & 1186 & +847 & \textit{critical} \\
\textbf{19} & 32.0 & 0.7 & 23.8 & 549 & +163 & \textit{critical} \\
\textbf{20} & 31.0 & 0.3 & 23.2 & 619 & -1487 & \textit{critical} \\
\textbf{21} & 28.6 & 0.2 & 20.9 & 801 & -1926 & \textit{critical} \\
\textbf{22} & 26.4 & 0.2 & 18.8 & 695 & -1441 & \textit{critical} \\
{23} & 24.6 & 0.2 & 17.2 & 536 & -1093 & \textit{dropped} \\
\bottomrule
\end{tabular}
\end{table}

\paragraph{Two validation checks.} (i) \emph{No cost-shock
confound}: CCGT price-setting share is $9.1\%$ critical vs $9.0\%$
flat --- gas-cost shocks cannot propagate asymmetrically across the
two classes, so within-day differencing identifies behaviour, not
cost. (ii) \emph{Outcome-side}: post-MTU15-DA the within-hour DA
quarter clearing-price SD averages $6.58$~EUR/MWh in critical hours
vs $2.13$ in flat ($3.1\times$), so the classifier correctly
predicts where post-reform quarter prices actually diverge.


% =====================================================================
\section{Who clears in which hours: the tech-mix by hour-class}\label{sec:prov:tech_mix}
% =====================================================================

\noindent Before turning to who price-sets, it is worth fixing the
basic descriptive picture: in each hour-class, what share of the
final cleared MWh actually comes from each technology? The cross-tab
in Table~\ref{tab:tech_mix_by_hourclass} is computed on the
\emph{final program} (PHF, post-IDA) for the post-DA15 window
(2025-10-01 to 2025-12-31), so it answers the natural question
``after every market has cleared, who is actually delivering the
energy in each hour-class?''

\begin{table}[H]
\centering\footnotesize
\setlength{\tabcolsep}{6pt}
\begin{tabular}{l r r r r}
\toprule
\textbf{Tech} & \textbf{Morning ramp (5--8)} & \textbf{Evening ramp (16--22)} & \textbf{Midday (11--14)} & \textbf{Flat (1--3)} \\
\midrule
Nuclear      & \textbf{34.3} & \textbf{26.9} & 25.8 & \textbf{34.0} \\
Wind         & 20.9 & 16.3 & 13.7 & 21.6 \\
\textbf{CCGT}         & \textbf{14.0} & \textbf{13.2} & 9.2 & \textbf{13.6} \\
\emph{Unknown}\,$^{\dagger}$ & 12.1 & 13.3 & 14.7 & 12.8 \\
Hydro        & 10.0 & 13.4 & 4.6 & 9.7 \\
\textbf{Cogen}        & 5.1 & 4.0 & 3.5 & 5.1 \\
Hydro\_RES   & 1.4 & 1.2 & 1.0 & 1.4 \\
Biomass      & 1.2 & 0.9 & 0.9 & 1.2 \\
Hydro\_pump  & 0.9 & 1.9 & 0.1 & 0.5 \\
Solar PV     & 0.0 & 8.7 & \textbf{26.4} & 0.0 \\
Coal         & 0.1 & 0.1 & 0.1 & 0.1 \\
\midrule
\textbf{Hour-class total (TWh)} & 13.9 & 31.1 & 18.5 & 10.5 \\
\bottomrule
\end{tabular}
\caption{\textbf{Final-program (PHF) share of cleared MWh by tech and hour-class, post-DA15 window.} Window 2025-10-01 to 2025-12-31 (92 days). Numbers are \% of total cleared MWh in each hour-class (column). PHF is the post-IDA program; we take the last-session assigned MW per (date, period, unit). The hour-class totals (TWh) reflect different hour counts per class per day (4 morning-ramp, 7 evening-ramp, 4 midday, 3 flat).
$^{\dagger}$``Unknown'' is the $\sim\!13\%$ of cleared MWh from unit codes not in the unit-map (mostly HISVD*/GALVD*/GNVD*/GSVD* virtual-delivery families --- portfolio/aggregator units that the unit-map drops); not a tech category.
Reproduction: \texttt{scripts/analysis/firm/tech\_mix\_by\_hour\_class.py}.}
\label{tab:tech_mix_by_hourclass}
\end{table}

\paragraph{Reading the cross-tab.}
\begin{itemize}\setlength{\itemsep}{2pt}
\item \textbf{Nuclear is the largest single bucket in every
non-midday hour-class} ($\sim\!34\%$ morning ramp and flat,
$\sim\!27\%$ evening). Inframarginal baseload with no within-hour
variation --- it sets the floor of the merit order in every hour.
\item \textbf{Solar PV} is the headline midday tech ($26.4\%$
midday, $\sim\!0\%$ overnight, $8.7\%$ in the evening ramp where the
ramp extends past sunset). The diurnal cycle is fully visible.
\item \textbf{Wind} is the second-largest in every hour-class except
midday ($\sim\!20\%$ morning and flat), dropping to $\sim\!14\%$
midday when solar takes share.
\item \textbf{CCGT} (the strategic-bidder pool of
\emph{preliminary results} \S 4.C) clears $\sim\!13$--$14\%$ of
total in every non-midday hour-class --- sizeable but not dominant.
Drops to $\sim\!9\%$ in midday when solar pushes it off the merit
order. The strategic-engagement story of \emph{preliminary results}
\S 4.C is therefore about restructuring the bid-shape of a
$\sim\!13\%$-of-cleared-MWh fleet, not about reshuffling the merit
order at any margin.
\item \textbf{Cogen} (separate from CCGT, see below) clears
$\sim\!3$--$5\%$ in every hour-class. \emph{Its share is roughly
flat across hour-classes}, consistent with the heat-driven (non-
strategic) dispatch story --- Cogen doesn't shift to follow the
power-side demand profile because its operating decision is driven
by industrial heat demand and a regulated remuneration regime, not
by the wholesale price.
\item \textbf{Hydro} is meaningfully higher in critical hours
(evening $13.4\%$, morning $10\%$) than in midday ($4.6\%$),
consistent with the reservoir-arbitrage story (dispatch into peaks).
\item \textbf{Hydro\_pump} is small everywhere but skews toward
critical evening ($1.9\%$ vs $0.1\%$ midday) --- discharges into the
evening peak, near-zero in midday when it would be pumping.
\end{itemize}

\paragraph{Why CCGT and Cogen are kept separate.} The unit-map distinguishes \texttt{CCGT} (58 large units, Big-4 dominated, two-tier strategic bid curve with Fase~I-recall economics) from \texttt{Cogen} (434 small units, regulated heat-driven dispatch under Real Decreto 413/2014 --- the ``$\sim\!5\%$ dispatchable tail'' of \emph{preliminary results} App.~C.1 footnote). Pooling would dilute the strategic-engagement story; the analysis throughout restricts to CCGT.

\subsection*{Cogen as a Spec~C null: a tech-scope robustness check}\label{sec:prov:cogen_stability}
\addcontentsline{toc}{subsection}{Cogen as a Spec~C null: a tech-scope robustness check}

\noindent Spec~C in \emph{preliminary results} \S 4.C runs on CCGT, hydro, pump-storage and wind. Cogen is excluded \emph{ex ante} on a regulated-remuneration prior: heat-driven dispatch should remove the within-hour optimisation problem that the granularity reform expands. We verify the prior by running the same Spec~C on Cogen at the same bandwidth ($h = 50$~EUR/MWh) and windows (pre 2025-04-28 to 2025-09-30, post 2025-10-01 to 2025-12-31); Biomass and Nuclear reported alongside.

\begin{table}[H]
\centering\footnotesize
\caption{\textbf{Spec~C robustness on technologies excluded from
\emph{preliminary results} \S 4.C.} Per-curve DiD of $\sigma_p$ and
$N_{\text{eff}}$ at the same DA15 DA bandwidth ($h = 50$~EUR/MWh)
and windows (pre 2025-04-28 to 2025-09-30; post 2025-10-01 to
2025-12-31) used in the headline memo. DiD = (post crit $-$ post
flat) $-$ (pre crit $-$ pre flat). CCGT row added for orientation
to the headline memo's reported value (\emph{preliminary results}
\S 4.C Table: CCGT DA $\sigma_p$ DiD $0.68$, $N_{\text{eff}}$ DiD
$1.89$). Per-cell sample sizes in the source CSV.
Reproduction: \texttt{scripts/analysis/firm/cogen\_stability\_check.py}.}
\label{tab:cogen_specc}
\setlength{\tabcolsep}{4pt}
\begin{tabular}{l r r r r r r}
\toprule
        & \multicolumn{2}{c}{$\sigma_p$ (EUR/MWh)} & & \multicolumn{2}{c}{$N_{\text{eff}}$} & \\
\cmidrule(lr){2-3} \cmidrule(lr){5-6}
\textbf{Tech} & post (crit, flat) & pre (crit, flat) & \textbf{DiD $\sigma_p$} & post (crit, flat) & pre (crit, flat) & \textbf{DiD $N_{\text{eff}}$} \\
\midrule
\textbf{CCGT} (in scope)   & 3.14,\ 2.18 & 2.55,\ 2.39 & \textbf{$0.79$} & 3.85,\ 2.40 & 2.85,\ 3.39 & \textbf{$1.99$} \\
\midrule
Cogen     & 0.71,\ 1.03 & 0.49,\ 0.74 & $-0.08$ & 1.04,\ 1.06 & 1.04,\ 1.05 & $-0.00$ \\
Biomass   & 1.74,\ 1.95 & 0.82,\ 0.69 & $-0.34$ & 1.57,\ 1.69 & 1.32,\ 1.35 & $-0.09$ \\
Nuclear   & 0.59,\ 0.55 & 1.83,\ 3.30 & $1.51$ & 1.65,\ 1.57 & 1.80,\ 1.74 & $0.03$ \\
\bottomrule
\end{tabular}
\end{table}

\textit{Reading.} Cogen $\sigma_p$ DiD $-0.08$, $N_{\text{eff}}$ DiD $\approx\!0$ --- the prior holds. Biomass behaves like Cogen (second regulated-remuneration tech, no granularity response). Nuclear's $1.51$ comes from a thin pre-flat cell ($n=233$, 4 units) and is not interpreted. CCGT reproduces \emph{preliminary results} \S 4.C within rounding ($0.79$ vs published $0.68^{***}$; $1.99$ vs $1.89^{***}$; unweighted-mean vs unit-FE coefficient).

\paragraph{Quantity-side mirror.} PHF cleared MWh/day stays flat for Cogen across all six regimes (Cogen 35.1 / 42.6 / 44.1 / 34.0 / 41.4 / 44.3 GWh/day for pre-IDA / 3-sess / ISP15-win / MTU15-IDA pre-blk / MTU15-IDA post-blk / DA15/ID15), against CCGT tripling ($39 \to 129$). Biomass similarly flat ($\sim 10$ GWh/day). Cogen operates under the regulated renewables-and-cogen remuneration framework (Real Decreto 413/2014): heat-driven dispatch, single-block must-clear bid at $\sim\!\texteuro 51$/MWh, no within-hour optimisation problem --- so granularity adds no instrument and produces no DiD signal.

\bigskip

% =====================================================================
\section{Who price-sets}\label{sec:prov:pricesetter}
% =====================================================================

\subsection{For the day-ahead market}\label{sec:prov:pricesetter_da}

\paragraph{Method.} The price-setter is the unit whose at-MCP step is \emph{partially} accepted: $q_{\text{below}} < q_{\text{assigned}} < q_{\text{below}} + q_{\text{at}}$ (with $q_{\text{at}}$ the MW bid at MCP $\pm 0.01$~EUR/MWh). Full acceptance (inframarginal) and curtailed-to-zero (minimum-income deactivation) excluded. Aggregations MW-weighted by $q_{\text{at}}$.

\paragraph{Who is marginal.} \textbf{Wind is the largest sell-side price-setter} in almost every cell ($18$--$71\%$ of at-MCP MW), reflecting renewable aggregators landing a single block at-the-money. Hydro and pump-storage next; in DA15/ID15, Iberdrola owns $\approx\!65\%$ critical and $80\%$ flat through pump, hydro and wind. \textbf{CCGT's share is small and volatile} ($1$--$17\%$ critical). Cross-check: CCGT critical/flat is near-symmetric (3-sess $8.0$/$8.9$, DA15/ID15 $2.7$/$2.0$), so a gas-cost shock --- which hits both classes alike --- is not what within-day differencing detects. Per-firm split within the CCGT pool:

\begin{table}[H]
\centering\scriptsize
\setlength{\tabcolsep}{5pt}
\begin{tabular}{l l r r r r r}
\toprule
Regime & Hour-class & IB & GE & GN & HC & OTHER \\
\midrule
3-sess & Critical & \textbf{60} & 31 & 0 & 0 & 8 \\
3-sess & Flat & 16 & \textbf{50} & 0 & 0 & 25 \\
ISP15-win & Critical & 35 & \textbf{57} & 0 & 0 & 5 \\
ISP15-win & Flat & 2 & \textbf{74} & 0 & 3 & 20 \\
DA60/ID15 post-blk & Critical & 31 & \textbf{69} & 0 & 0 & 0 \\
DA60/ID15 post-blk & Flat & 19 & \textbf{80} & 0 & 0 & 0 \\
\textbf{DA15/ID15} & Critical & \textbf{65} & 29 & 0 & 6 & 0 \\
\textbf{DA15/ID15} & Flat & \textbf{80} & 0 & 6 & 14 & 0 \\
\bottomrule
\end{tabular}
\caption{\textbf{Per-firm shares within the CCGT price-setter pool} (\% of CCGT $q_{\text{at}}$). \textbf{GN's CCGT bid stack is so heavily concentrated far above the clearing-price ceiling (a flat shelf at $\sim\!1{,}000$ EUR/MWh) that GN essentially never lands at MCP}. IB and GE alternate as the dominant CCGT price-setter. The 33.1\% CCGT flat-hour figure in DA60/ID15 post-blackout is 80\% GE.}
\label{tab:ccgt_by_firm}
\end{table}

\paragraph{Price-setting $\neq$ market power for CCGT.} A unit is flagged ``marginal'' merely for having offered MW at the price that happened to clear. The ranking actually \emph{inverts} the market-power ordering: GN withholds the most ($75\%$ parked, \S\ref{sec:prov:per_curve_def}) but scores $0\%$ in the price-setter pool precisely \emph{because} it bids too high ever to be marginal; IB tops the pool only by keeping a thin competing tier near MCP. CCGT market power is exercised by withholding \emph{from} the day-ahead and captured downstream in pre-IDA Fase~I redispatch (\S\ref{sec:prov:phf_minus_pdbf}); we read CCGT conduct from bid-vs-marginal-cost (\S\ref{sec:prov:per_curve_def}) and the Fase~I channel, not from the price-setter share.

\subsection{For the intraday markets}\label{sec:prov:pricesetter_ida}

\noindent The same partial-acceptance rule, applied to the intraday-auction clearing prices, bid stack and programs (sell side, pooled across the 3 IDA sessions per day), shows \textbf{CCGT far more active as an IDA price-setter than as a DA one}: $24.2\%$ of critical-hour IDA mass under DA15/ID15 against $2.7\%$ in DA --- a factor of $\sim 9$. This is the price-setter-side trace of the day-ahead-to-intraday substitution documented in \S\ref{sec:prov:per_curve_def}: CCGTs bid above the day-ahead clearing price but compete in the intraday auction, where they clear at the marginal step. Wind still dominates ($\sim 32\%$ critical) but by a smaller margin than in DA.


% =====================================================================
\section{Bid-curve shape: $\sigma_p$ and $N_{\text{eff}}$}\label{sec:prov:dw_decomp}
% =====================================================================

\noindent We use the same two per-curve bid-shape scalars as the headline memo: $\sigma_p$, the MW-weighted SD of in-band tranche prices, and $N_{\text{eff}}$, the inverse Herfindahl of in-band tranche MW shares. Both are defined per single bid curve --- one value per (unit, date, period) sell curve --- with no comparison between curves. Cross-regime, cross-firm and cross-hour reading is then a comparison of the \emph{values} $\sigma_p$ and $N_{\text{eff}}$ across curves.

% =====================================================================
\subsection{Definitions and bandwidth}\label{sec:prov:per_curve_def}
% =====================================================================

\paragraph{Bandwidth $h$: same as the headline memo.} Both metrics restrict to a price band $|p - \text{MCP}| \le h$ --- the region where a bid is in contention. \textbf{The bandwidth is the headline memo's} (\emph{preliminary results} \S 1, Table of bandwidths): the window-and-market-specific $p_{90} - p_{50}$ of the MCP distribution in each (reform, market) cell, giving $h \in \{45, 46, 49, 50, 58, 62\}$~EUR/MWh across the eight (reform, real-or-placebo, DA-or-IDA) cells of the headline memo. This is the \emph{strategically active competing core}: the price region in which the bid actually competes for clearance.

The MW-weighted density of $|p_{\text{bid}} - \text{MCP}|$ over all day-ahead sell tranches of the price-setting techs is \textbf{bimodal}: a dense near-MCP competing cluster (where the window-specific $p_{90}$ bands sit) and a sparse far scarcity-cap cluster (the parked withholding tier). The two clusters meet at a sharp cliff around $\sim\!140$~EUR/MWh, which coincides with the realised clearing-price ceiling (median MCP $+\,140 \approx p_{99.9}$ of MCP). The per-curve metrics table in \S\ref{sec:prov:dw_tables} at a single pooled $h = 140$ is reported as a \emph{wider-band descriptive supplement}: it describes the full competitive distribution out to the cliff, useful for tech-comparison context but not for the Spec~C identification. The identification reads in this document (\S\ref{sec:prov:cogen_stability}, \S\ref{sec:prov:per_firm_tech_shape}) all use the headline memo's window-and-market-specific $h$.

\paragraph{The two-tier bid curve, read against marginal cost.} The bimodality of Figure~\ref{fig:strategic_band} is, unit by unit, a \textbf{two-tier bid curve}: a \emph{competing tier} around marginal cost (clears the auction) plus a \emph{scarcity tier} parked far above any plausible clearing price (does not clear --- deliberate withholding). The parked capacity is recalled in the pre-IDA Fase~I redispatch through a \emph{separate} restriction offer (pay-as-bid), so the day-ahead scarcity-tier price is a non-payment focal choice.\footnote{P.O.~3.2 process; settlement under P.O.~14.4 apartado~20 (\texttt{BOE-A-2025-5342}, \texttt{BOE-A-2025-13076}). Bidding high in DA to be excluded and then dispatched in restrictions is the conduct CNMC sanctioned in 2019 (\texttt{SNC/DE/174/17}, \texttt{SNC/DE/175/17}).} The withholding boundary is the $\sim\!200$~EUR/MWh clearing-price ceiling. Aggregated over DA15/ID15, the withheld share above the ceiling is firm-uniform in spirit but very different in magnitude: Naturgy $75\%$ (flat shelf at $\sim\!1{,}000$~EUR/MWh), Endesa $22\%$ (much at the $4{,}000$ cap), Iberdrola $22\%$ (a $200$--$400$ ramp), EDP-HC $19\%$. Iberdrola looks competitive against a $500$~EUR/MWh cutoff but withholds as much as Endesa against the true $\sim\!200$ ceiling, and clears only $\sim\!16\%$ of its final programme in the day-ahead (barely above Naturgy's $\sim\!14\%$). The firms differ in \emph{how far} they bid above the ceiling, not in \emph{whether} they withhold. The intraday auction confirms it is strategy, not cost: on the same units and days, all four fleets bid a single competitive tier at MW-weighted $82$--$96$~EUR/MWh with $\sim\!0\%$ withheld --- Naturgy at $90$~EUR/MWh against its $767$ day-ahead. The per-curve functionals below restrict to the competing tier, where price-graded bidding is the choice variable.

\paragraph{Per-curve functionals and offer-type caveat.} $\sigma_p$ (MW-weighted SD of in-band tranche prices) and $N_{\text{eff}}$ (inverse Herfindahl of in-band tranche MW shares) are defined in \emph{preliminary results} \S 1; we use the same definitions here. Block orders mechanically force $\sigma_p\!=\!0$ and $N_{\text{eff}}\!=\!1$ and MIC orders displace shaping onto the income floor, so the offer-type mix matters: wind, solar, hydro and pump-storage are $\sim\!100\%$ simple; nuclear $64\%$ simple; \textbf{CCGT is only $18\%$ simple ($50\%$ MIC, $32\%$ block)}. For CCGT, any geometry-to-conduct step conditions on offer type (\emph{preliminary results} Appendix A.7 reports the type-split DiD).

% =====================================================================
\subsection{Results}\label{sec:prov:dw_tables}

\noindent The results below report $\sigma_p$ and $N_{\text{eff}}$ at a wider supplementary bandwidth ($h = 140$) as cross-regime descriptive context. The Spec~C identification at the headline memo's window-and-market-specific bandwidth lives in \emph{preliminary results} \S 4.C (per-tech pooled DiD), \S\ref{sec:prov:per_firm_tech_shape} (per-firm $\times$ tech), and \S\ref{sec:prov:cogen_stability} (Cogen / Biomass / Nuclear robustness).

\paragraph{Per-curve functionals at a wider supplementary bandwidth ($h = 140$).}
Table~\ref{tab:per_curve_metrics} and Figure~\ref{fig:per_curve_metrics} report $\sigma_p$ and $N_{\text{eff}}$ averaged over (unit, date, period) day-ahead sell-bid curves, by tech, hour-class and regime, at a single pooled $h = 140$~EUR/MWh --- the empirical cliff between the competing cluster and the parked tier (Figure~\ref{fig:strategic_band}). \emph{This is a wider-band supplementary view, not the Spec~C identification.} The Spec~C identification uses the headline memo's window-and-market-specific bandwidth (\S\ref{sec:prov:per_curve_def}) and is read in \S\ref{sec:prov:cogen_stability} (Cogen Spec~C null) and \S\ref{sec:prov:per_firm_tech_shape} (per-(firm, tech) Spec~C DiD); the headline DiD on the dispatchable techs lives in \emph{preliminary results} \S 4.C. The wider band here adds tech-comparison context (CCGT's competing-cluster tail extends further from MCP than hydro's, for instance) at the cost of mixing tranches that would not be picked up by the tighter $p_{90}$ band.

\begin{table}[H]
\centering\footnotesize
\caption{\textbf{Per-curve bid-shape metrics by tech, hour-class and regime (wider-band supplementary, $h = 140$).} DA sell bids, in-band $|p-\text{MCP}|\le 140$~EUR/MWh (the descriptive cliff bandwidth, not the headline memo's $p_{90}-p_{50}$ bandwidth used for Spec~C identification). Mean over (unit, date, period) curves. $\sigma_p$ = MW-weighted price spread (EUR/MWh); $N_{\text{eff}}$ = effective in-band tranche count. $953{,}384$ in-band curves. Cross-regime reading: tech ranking holds, but the absolute magnitudes are inflated relative to the Spec~C bandwidth by the inclusion of the outer tail of the competing cluster.}
\label{tab:per_curve_metrics}
\setlength{\tabcolsep}{4pt}
\begin{tabular}{l l r r r r r r r r r r }
\toprule
 & & \multicolumn{2}{c}{3-sess} & \multicolumn{2}{c}{ISP15-win} & \multicolumn{2}{c}{DA60/ID15 pre} & \multicolumn{2}{c}{DA60/ID15 post} & \multicolumn{2}{c}{DA15/ID15} \\
Tech & Hour & $\sigma_p$ & $N_{\text{eff}}$ & $\sigma_p$ & $N_{\text{eff}}$ & $\sigma_p$ & $N_{\text{eff}}$ & $\sigma_p$ & $N_{\text{eff}}$ & $\sigma_p$ & $N_{\text{eff}}$ \\
\midrule
CCGT & Critical & 36.1 & 3.3 & 32.0 & 4.2 & 29.2 & 3.1 & 24.5 & 3.1 & 25.5 & 3.4 \\
CCGT & Flat & 27.5 & 2.5 & 45.2 & 2.7 & 17.1 & 2.3 & 20.2 & 2.6 & 18.4 & 2.2 \\
\addlinespace
Hydro & Critical & 17.5 & 4.0 & 16.3 & 4.0 & 13.1 & 3.5 & 15.5 & 3.5 & 18.3 & 3.9 \\
Hydro & Flat & 17.6 & 4.0 & 18.5 & 3.8 & 12.6 & 3.5 & 16.0 & 3.5 & 17.0 & 3.9 \\
\addlinespace
Hydro pump & Critical & 3.6 & 2.5 & 3.8 & 2.5 & 2.6 & 2.3 & 3.9 & 2.8 & 4.3 & 3.5 \\
Hydro pump & Flat & 2.1 & 2.6 & 3.9 & 2.9 & 2.4 & 2.3 & 2.6 & 2.4 & 2.9 & 4.2 \\
\addlinespace
\bottomrule
\end{tabular}
\end{table}

\begin{figure}[H]
\centering
\includegraphics[width=0.82\textwidth]{../../figures/working/per_curve_metrics.pdf}
\caption{\textbf{Per-curve bid-shape metrics across regimes (wider-band supplementary, $h = 140$).} Same caveat as Table~\ref{tab:per_curve_metrics}: wider-band descriptive view, not the Spec~C identification (Spec~C uses the window-specific $p_{90}-p_{50}$ bandwidth of the headline memo). Left: price spread $\sigma_p$; right: effective tranche count $N_{\text{eff}}$. Each line is a (tech, hour-class) cell; solid = critical hours, dashed = flat. CCGT carries the widest price spread ($\sigma_p \approx 25$--$40$ EUR/MWh at this wider band) and it is larger in critical than flat hours in four of the five regimes --- the qualitative pattern that Spec~C identifies; Hydro is intermediate ($\sim 13$--$18$) and Hydro pump near-flat ($\sim 3$--$4$).}
\label{fig:per_curve_metrics}
\end{figure}

\paragraph{Wind within-hour variation is forecast, not strategy.} Regressing each wind curve's $\sigma_p$ on a within-hour scarcity proxy (residual demand net of wind, deviated from clock-hour mean) with unit$\times$date$\times$hour FE (DA15/ID15, $2.3$M in-band wind curves, date-clustered SEs): $84\%$ of wind in-band curves are a single block and $89\%$ of clock-hours show no across-quarter $\sigma_p$ variation; where $\sigma_p$ does move, the scarcity response is $0.01$~EUR/MWh per GW (within-hour partial correlation $0.01$) --- economically negligible. The largest aggregator response, Gesternova's $\beta\!=\!0.14$, is still tiny against $\beta\!\approx\!0$ for AXPO, NEXUS, IGNIS. Within-hour wind movement is the sub-hourly forecast, not strategic scarcity-chasing; wind reads as a price-taking baseline.

\begin{figure}[H]
\centering
\includegraphics[width=0.80\textwidth]{../../figures/working/wind_within_hour_scarcity.pdf}
\caption{\textbf{The wind functional price metric $\sigma_p$ does not respond to within-hour scarcity} (DA15/ID15, unit$\times$date$\times$hour fixed effects; band centred on the clock-hour-mean clearing price). Left: binned within-hour-demeaned $\sigma_p$ against the within-hour-demeaned scarcity proxy (load minus solar forecast, GW), for critical and flat hours --- both slopes essentially flat. Right: the slope for each clock-hour of the day (red = critical, blue = flat hours; date-clustered $95\%$ intervals) --- no hour reaches an economically meaningful response.}
\label{fig:wind_within_hour}
\end{figure}

\paragraph{Pump-storage within-day arbitrage under MTU15-DA: a descriptive look (no causal claim).} The bid-shape evidence above is one side of the granular-pricing reform; one would naturally look at pump-storage for the quantity-side mirror. Pre/post means on the two sides of the pump-storage book, taken at face value, show a striking within-day reallocation:

\begin{center}
\footnotesize
\setlength{\tabcolsep}{8pt}
\begin{tabular}{l r r c r r}
\toprule
& \multicolumn{2}{c}{\textbf{Generation (MW per period)}} & & \multicolumn{2}{c}{\textbf{Pumping (MW per period)}} \\
\cmidrule(lr){2-3}\cmidrule(lr){5-6}
\textbf{Hour-class} & Pre (Jul--Sep) & Post (Oct--Dec) & & Pre & Post \\
\midrule
Critical & 851 & 1\,018 ($20\%$) & & 664 & 317 ($-52\%$) \\
Flat     & 724 & 258 ($-64\%$)    & &  52 & 566 ($10\times$) \\
\bottomrule
\end{tabular}
\end{center}

\noindent Generation appears to reallocate from flat hours
($-64\%$) to critical hours ($20\%$); pumping mirrors it. Total
daily generation is essentially unchanged ($14\,874 \to 14\,250$
MWh, $-4\%$). The reading would be textbook within-day arbitrage:
pre-MTU15-DA, hourly bids could not price-discriminate across
quarters; post-reform, the unit class with the largest within-hour
arbitrage value exploits within-hour price spreads cleanly.

\emph{The headline-methodology verdict is more cautious.} The BSTS
counterfactual on per-tech cleared MWh (Spec~A in
\emph{preliminary results} \S 4.A) does not retain the pump-storage
MTU15-DA effect under the extrapolation-drift diagnostic (see
\emph{preliminary results} \S 4.A.6, ``BSTS extrapolation drift
cases''). The pre/post raw means above remain descriptively
informative --- the within-day reallocation pattern is real in the
data --- but we do not interpret them as a causal MTU15-DA effect
for the reasons set out in the preliminary memo (cross-season pre/
post, drift between the 14-day and 40-day post-windows, sensitivity
to the choice of pre-window).

\subsection{XBID extension of Spec~C}\label{sec:prov:xbid_specc}

\noindent Third-venue test of preliminary's Spec~C: $\sigma_p$ and $N_{\text{eff}}$ on Big-4 CCGT XBID sell orders (\texttt{exec\_condition} $\in\{$\,'', ICE$\}$), in-band ($h = 62$~EUR/MWh) around the same-period IDA-auction clearing as anchor, around the ID15 cutover. Pre 2024-12-19 to 2025-03-18 ($\sim\!90$ days), post 2025-03-19 to 2025-04-27 (40 days, pre-blackout).

\begin{center}\footnotesize
\setlength{\tabcolsep}{6pt}
\begin{tabular}{l l r r r r}
\toprule
Window & Hour-class & $n_{\text{curves}}$ & $\overline{n_{\text{tr}}}$ & $\sigma_p$ & $N_{\text{eff}}$ \\
\midrule
pre  ID15 & Critical & 3{,}710 & 16.2 & 3.45 & 15.10 \\
pre  ID15 & Flat     & 1{,}366 & 12.9 & 3.55 & 12.27 \\
post ID15 & Critical & 2{,}874 & 13.1 & 2.81 & 12.23 \\
post ID15 & Flat     &   683 & 13.3 & 2.11 & 13.14 \\
\midrule
\multicolumn{5}{r}{\textbf{DiD $\sigma_p$:}} & \textbf{$0.81$} \\
\multicolumn{5}{r}{\textbf{DiD $N_{\text{eff}}$:}} & \textbf{$-3.73$} \\
\bottomrule
\end{tabular}
\end{center}

\textit{Reading.} Critical-vs-flat CCGT XBID $\sigma_p$ \emph{widens} by $0.81$~EUR/MWh post-ID15 (same sign as the auction-side hydro/pump DiDs in \emph{preliminary results} \S 4.C). $N_{\text{eff}}$ \emph{tightens} by $-3.73$ effective tranches: with the IDA auction itself moving to 15-min, CCGT XBID orders become less granular in critical hours --- the continuous-market shaping load shifts to the now-finer IDA auction. The strategic-latitude story extends to a third venue, redistributively. Reproduction: \texttt{scripts/analysis/bid/xbid\_specc\_did.py}.

\subsection{Per-firm $\times$ technology Spec~C decomposition}\label{sec:prov:per_firm_tech_shape}

\noindent The pooled per-(tech, market) DiD coefficients in
\emph{preliminary results} \S 4.C summarise the bid-shape response
to MTU15-DA at the technology level. The per-firm CCGT split
appears once, in \emph{preliminary results} \S A.4 (``primarily
Naturgy, Iberdrola to a lesser extent, Endesa and EDP-HC do not
contribute''). This subsection extends the per-firm split to all
dispatchable techs in the Spec~C scope (CCGT, hydro, pump-storage,
wind) at the \emph{same} per-curve outcomes ($\sigma_p$,
$N_{\text{eff}}$), the \emph{same} window-and-market-specific
bandwidth (DA15 DA $h = 50$~EUR/MWh, DA15 IDA $h = 58$~EUR/MWh),
and the \emph{same} pre/post windows used in the headline memo.
This is per-firm Spec~C, not a new spec.

\begin{table}[H]
\centering\footnotesize
\caption{\textbf{Per-(firm, tech) DA15 DA Spec~C DiD on $\sigma_p$ and
$N_{\text{eff}}$, $h = 50$~EUR/MWh.} DiD $= (\sigma_p$ post crit
$- $ post flat$) - ($pre crit $-$ pre flat$)$, computed cell-by-cell
on the per-curve panel of \emph{preliminary results} \S 4.C
(pre 2025-07-01 to 2025-09-30; post 2025-10-01 to panel end
2025-11-09). Sign and magnitude reproduce \emph{preliminary results}
\S A.4 attribution for CCGT (GN primary, IB lesser, GE near zero,
HC small). Per-tech leader differs: GN on CCGT, IB and HC on hydro,
GE on pump. Empty cells: firm has no units of that tech in the
Spec~C panel. Reproduction:
\texttt{scripts/analysis/firm/per\_firm\_tech\_bid\_shape.py}.}
\label{tab:per_firm_tech_specc}
\setlength{\tabcolsep}{5pt}
\begin{tabular}{l r r r r | r r r r}
\toprule
        & \multicolumn{4}{c}{\textbf{DiD $\sigma_p$ (EUR/MWh)}} & \multicolumn{4}{c}{\textbf{DiD $N_{\text{eff}}$}} \\
\cmidrule(lr){2-5} \cmidrule(lr){6-9}
\textbf{Firm} & CCGT & Hydro & Pump & Wind & CCGT & Hydro & Pump & Wind \\
\midrule
IB    & $0.96$         & \textbf{$2.44$} & $0.24$         & $-0.06$ & $0.21$ & $0.51$ & $-0.12$ & $-0.03$ \\
GN    & \textbf{$3.72$} & $1.39$         & $0.17$         & $0.01$ & \textbf{$7.91$} & $0.06$ & $0.09$ & $-0.08$ \\
GE    & $-0.10$         & $0.81$         & \textbf{$0.85$} & ---     & $0.01$ & $-0.01$ & $0.01$ & ---     \\
HC    & $0.83$         & \textbf{$2.18$} & ---             & $-0.27$ & $0.34$ & $0.74$ & ---     & $-0.06$ \\
\midrule
\textbf{Headline} (\emph{prelim.} \S 4.C) & \textbf{$0.68^{***}$} & \textbf{$1.80^{***}$} & \textbf{$0.92^{***}$} & \textbf{$-0.02$} & \textbf{$1.89^{***}$} & \textbf{$0.51^{***}$} & $-0.06$ & $-0.01$ \\
\bottomrule
\end{tabular}
\end{table}

\paragraph{Reading.} (i) \emph{The CCGT per-firm pattern matches
\emph{preliminary results} \S A.4.} GN's $\sigma_p$ DiD is $3.72$
(the largest), IB's is $0.96$, HC's $0.83$, GE's near zero. The
$N_{\text{eff}}$ DiD is even more concentrated on GN ($7.91$
effective tranches, vs $0.21$ for IB) --- GN's competing tier
restructures into much finer per-quarter ladders, GN does most of
the work in the pooled CCGT $1.89$ $N_{\text{eff}}$ headline. (ii)
\emph{The per-tech leader differs by tech.} GN leads on CCGT, IB
and HC on hydro, GE on pump-storage. A firm that engages with the
granularity instrument on one technology does not automatically
engage on every technology in its portfolio --- GN's hydro and pump
$\sigma_p$ DiDs are smaller than IB's and HC's, and GN's
$N_{\text{eff}}$ on hydro and pump is essentially zero. (iii)
\emph{Wind is uniformly null} at the per-firm level too, consistent
with the operational-uniform-across-hour-classes reading of
\emph{preliminary results} \S 4.C.2 (the engagement is on the
volume side, not the bid-shape side). (iv) \emph{The pooled
headline rows} reproduce \emph{preliminary results} \S 4.C within
rounding (unweighted-mean vs unit-FE regression coefficient
explains the $\sim 0.1$--$0.2$ gap on each pooled coefficient).

\paragraph{Implication.} Each pooled per-tech DiD of \emph{preliminary results} \S 4.C is a portfolio-weighted sum of a small number of (firm, tech) cells: the CCGT $\sigma_p$ widening is structurally a GN-CCGT event with small IB support; the hydro widening is IB-and-HC; the pump widening leans on GE. Each dispatchable firm responds on the technology where their portfolio gives them within-hour optimisation room.

% =====================================================================
\section{Continuous-intraday repositioning (short)}\label{sec:prov:ci_per_tech}
% =====================================================================

\noindent XBID volume tracks the granular-auction story: gross XBID activity is $\sim\!32 \to 12 \to 23$~GWh/day across the five regimes ($-29\%$ baseline to DA15/ID15, trough $-62\%$ in the 40-day MTU15-IDA-only window). \textbf{CCGT is a systematic net seller} in every regime ($0.8$ to $2.5$~GWh/day) --- the continuous-market face of ``stay out of the day-ahead''. Within-hour rebalancing intensity (cross-quarter SD of XBID sell volume) is $2$--$5\times$ larger in critical than flat hours, ranked IB $>$ GN $>$ HC $>$ GE, and stable across post-MTU15-IDA regimes: CI \emph{volume} fell with the granular auction; per-firm rebalancing \emph{intensity} did not.


% =============================================================================
\part{Part B --- Regulatory and confounder context}
% =============================================================================

\paragraph{REE program-and-intervention chain (brief primer).} Each delivery day's unit program is built up as PDBC (DA-cleared) $\to$ PDBF ($+$ bilaterals) $\to$ PDVP ($+$ pre-IDA \emph{Fase I/II} RT, pay-as-bid) $\to$ PHF ($+$ IDA $+$ post-IDA RT2) $\to$ PHFC ($+$ continuous) $\to$ P48 ($+$ real-time TR). \emph{Reforzada} = REE's post-blackout (2025-04-28 onward) stance of larger pre- and post-IDA RT, mostly Fase~I-up. Roadmap: $q_1$ drop (\S\ref{sec:prov:q1_drop}), absorbing mechanism (\S\ref{sec:prov:phf_minus_pdbf}), zonal CCGT concentration (\S\ref{sec:prov:geo}), system-cost view (\S\ref{sec:prov:full_cascade}), CNMC enforcement (\S\ref{sec:prov:cournot_cnmc}).

% =====================================================================
\section{The $q_1$ drop: CCGT auction-cleared collapse}\label{sec:prov:q1_drop}
% =====================================================================

\noindent $q_1$ = sum of PDBC (OMIE DA auction-cleared) MWh, restricted to Big-4 (IB, GE, GN, HC), per (technology, calendar month); same-calendar-month comparison absorbs seasonality by construction. $q_{\text{final}} = $ PDBC $+$ Bilaterals $+$ (IDA $+$ REE-RT). The Big-4 aggregate by tech appears in \emph{preliminary results} Appendix~C.4 and is not reproduced; the figure below adds the per-CCGT-firm split, showing the DA $\to$ RT shift is firm-wide.

\begin{figure}[H]
\centering
\includegraphics[width=0.82\textwidth]{../../figures/working/fig_q_components_ccgt_by_firm_weekly.pdf}
\caption{\textbf{Weekly three-component decomposition per CCGT firm.} PDBC (blue, DA-cleared) $+$ bilateral executions (green) $+$ IDA $+$ REE real-time top-up (red) $=$ final post-IDA dispatch (black), Big-4 CCGT firms split. The DA~$\to$~RT shift is firm-wide and roughly tracks zonal CCGT dominance. Vertical guides: MTU15-IDA (2025-03-19), Blackout (2025-04-28), MTU15-DA (2025-10-01).}
\end{figure}

\begin{table}[H]
\centering
\caption{\textbf{Same-calendar-month auction-cleared DA volume by technology, GWh} (PDBC, pivotal-firm aggregate). Bilaterals excluded. Each column pair compares the same calendar month one year apart, so seasonality cancels. Only CCGT collapses; every other tech grows or stays flat.}
\label{tab:q1_drop_by_tech}
\begin{tabular}{l r r r r r r r r r}
\toprule
 & \multicolumn{3}{c}{Dec 2024 $\to$ Dec 2025} & \multicolumn{3}{c}{Nov 2024 $\to$ Nov 2025} & \multicolumn{3}{c}{Aug 2024 $\to$ Aug 2025} \\
\cmidrule(lr){2-4}\cmidrule(lr){5-7}\cmidrule(lr){8-10}
Tech & Pre & Post & \% & Pre & Post & \% & Pre & Post & \% \\
\midrule
\textbf{CCGT} & \textbf{1444} & \textbf{244} & \textbf{-83\%} & \textbf{637} & \textbf{130} & \textbf{-80\%} & \textbf{532} & \textbf{360} & \textbf{-32\%} \\
Nuclear & 1253 & 2587 & +106\% & 933 & 1909 & +105\% & 1997 & 2243 & +12\% \\
Hydro & 1016 & 640 & -37\% & 509 & 611 & +20\% & 303 & 424 & +40\% \\
Hydro\_pump & 195 & 247 & +27\% & 182 & 341 & +87\% & 239 & 406 & +70\% \\
Wind & 732 & 696 & -5\% & 666 & 1081 & +62\% & 379 & 449 & +18\% \\
Solar PV & 57 & 34 & -40\% & 37 & 100 & +170\% & 37 & 136 & +268\% \\
\bottomrule
\end{tabular}
\end{table}

\textit{Reading.} CCGT Big-4: Dec 2024 = 1{,}444~GWh $\to$ Dec 2025 = 244~GWh (\textbf{$-83\%$}); Nov $-80\%$. Every other tech grows or stays flat. The drop is the extensive margin (number of unit-hours cleared), not the intensive margin (per-cell mean MWh, stable at $\approx\!300$). Firm-wide critical-hour drops Dec24$\to$25: IB $-11\%$, GE $-6\%$, GN $-10\%$, HC $-18\%$ (flat-hour drops larger: GN $-30\%$, HC $-50\%$). Only CCGT can stay out of DA and rejoin via Fase~I and the intraday auctions --- renewables are not called up in PO~3.2 (energy-constrained), nuclear is must-clear. CCGT firms with zonal post-clearing backstop (\S\ref{sec:prov:geo}) bid above DA, are rescued by pay-as-bid Fase~I, and post a near-constant cap-padded DA ladder --- the practice CNMC sanctioned in 2019, now at fleet-wide scale.

\textit{Reconciliation with the headline memo.} \emph{Preliminary results} \S 4.A.2 reports a CCGT DA-cleared \emph{scale-up} of $40$~GWh/day under MTU15-DA. The BSTS scale-up is the DA15 increment \emph{conditional on a reforzada-era pre-window} (train 2025-04-28 to 2025-09-30, both pre and post post-reforzada); the $-83\%$ here is the \emph{joint} reform-and-reforzada effect against pre-reforzada Dec 2024. Together: the reforzada onset slashed CCGT DA dramatically; within the reforzada era, DA15 partially restored it in critical hours.

\subsection{Robustness: bilateral channel and outages}\label{sec:prov:q1_drop:bilat}

\noindent \textbf{Bilateral} (PDBF $-$ PDBC). Could CCGT firms have relabelled output as bilateral? Table~\ref{tab:bilat_check}: for CCGT the channel is small ($\sim\!90$--$128$~GWh/month, no upward trend), so the auction-cleared collapse is real.

\begin{table}[H]
\centering
\caption{\textbf{Bilateral channel by technology and regime, GWh / month} (PDBF $-$ PDBC, pivotal-firm aggregate). CCGT bilaterals are small and stable; Hydro / Nuclear / Wind bilaterals are an order of magnitude larger and structurally part of their dispatch.}
\label{tab:bilat_check}
\setlength{\tabcolsep}{4pt}
% auto-built by scripts/analysis/firm/bilat_channel_check.py
% Bilateral channel (PDBF - PDBC), GWh/month mean per regime
\begin{tabular}{l r r r r r}
\toprule
 & 3-sess & ISP15-win & DA60/ID15 pre-blk & DA60/ID15 post-blk & DA15/ID15 \\
\midrule
CCGT & 90 & 117 & 0 & 46 & 128 \\
Hydro & 2273 & 3627 & 4099 & 2537 & 3348 \\
Hydro pump & 5 & 15 & 61 & 37 & 40 \\
Nuclear & 2780 & 2887 & 638 & 2319 & 2449 \\
Wind & 716 & 605 & 371 & 642 & 1464 \\
Solar PV & 373 & 216 & 412 & 547 & 302 \\
\bottomrule
\end{tabular}

\end{table}

\textbf{Outages.} Aggregating ENTSO-E planned + forced outage windows (PSR B04) per regime, summing $(\text{nominal}_{\text{MW}} - \text{min\_avail}_{\text{MW}})$ per outage, against the $\sim\!23.6$~GW CCGT fleet (50 of 58 units covered):

\begin{center}\scriptsize
\setlength{\tabcolsep}{4pt}
\begin{tabular}{l r r r r}
\toprule
Regime & unavail MW (forced, planned) & avail. MW \\
\midrule
pre-IDA            & 4\,741 (815, 3\,925)        & 18\,869 \\
3-sess             & 3\,962 (762, 3\,200)        & 19\,648 \\
ISP15-win          & 5\,538 (683, 4\,855)        & 18\,072 \\
MTU15-IDA pre-blk  & 5\,692 (683, 5\,009)        & 17\,918 \\
MTU15-IDA post-blk & 4\,895 (726, 4\,170)        & 18\,715 \\
\textbf{DA15/ID15} & 4\,885 (\textbf{2\,389}, 2\,495) & \textbf{18\,725} \\
\bottomrule
\end{tabular}
\end{center}

\noindent CCGT available capacity is essentially flat across regimes ($\sim\!18.7$--$19.6$~GW). The $-83\%$ Dec24 $\to$ Dec25 q$_1$ drop \emph{cannot} be explained by reduced fleet availability. Forced outages tripled in DA15/ID15 ($726 \to 2{,}389$~MW) but planned outages dropped by an equal amount. ENTSO-E A73 metered output per CCGT unit (53 of 58, $\sim\!91\%$ coverage) gives $\sim\!49$~GWh/day physical generation post-DA15 against PHF-scheduled $\sim\!78$~GWh/day --- CCGT under-delivers the post-IDA program by $\sim\!37\%$ via real-time TR (downward). Reproduction: \texttt{scripts/analysis/regulatory/ccgt\_outage\_controls.py}.

% =====================================================================
\section{REE post-clearing inclusion as the absorbing mechanism}\label{sec:prov:phf_minus_pdbf}
% =====================================================================

\subsection{Full post-clearing intervention by \emph{tipo redespacho} code}\label{sec:prov:phf_minus_pdbf:fullview}

\begin{figure}[H]
\centering
\includegraphics[width=0.85\textwidth]{../../figures/working/fig_ree_intervention_full.pdf}
\caption{\textbf{Full REE post-clearing intervention by \emph{tipo redespacho} code, ESIOS archive 28 (\texttt{totalrp48preccierre}).} Top: UP-redispatch GWh / month; bottom: DOWN-redispatch (negative axis). Labels follow the project-canonical reference \S13: 33 $=$ real-time TR (general); 34 $=$ inter-zonal / network TR; \textbf{61 $=$ system-security TR under PO 3.2 (the post-IDA / real-time channel)}; 68 $=$ reserve management; 69 $=$ voltage control / black-start; 81 $=$ catch-all ``other''; 92 $=$ mFRR activation; 94 $=$ system balancing. Codes not documented in the project reference (19, 22, 24, 32, 38, 65, 66, 80, 82, 85, 89, 95, 96) are pooled into \emph{Uncategorised}. Code 23 ($\sim$10.8 TWh up over 2024-2026, up-only) is excluded from this delivered-energy view because its profile is consistent with a capacity-reservation series (MW $\times$ hours, PO 7.2 banda-secundaria); stacking it on the same axis as activated-energy channels would be apples-to-oranges.}
\label{fig:ree_intervention_full}
\end{figure}

\subsection{Per-firm $\times$ per-technology Path~A reading}\label{sec:prov:phf_minus_pdbf:perfirm}

\noindent \textbf{Programs: PHF $-$ PDBF, per (firm, technology, month).} Pre window 2024-01 $\to$ 2025-03; post window 2025-05 $\to$ 2026-01 (\emph{operación reforzada}). Bilaterals cancel out (in both terms at the same MW); the residual is IDA $+$ REE post-IDA RT.

\begin{table}[H]
\centering
\caption{\textbf{Post-DA intervention by firm and technology, GWh / month.} \textbf{Programs: PHF $-$ PDBF}, same construction as Table~\ref{tab:gap_pretrends} but resolved per firm. The bilateral channel (PDBF $-$ PDBC) is small for CCGT (Table~\ref{tab:bilat_check}); does \emph{not} drive this gap.}
\label{tab:phf_pdbf_alltech}
\begin{tabular}{l l r r r r}
\toprule
Tech & Firm & Pre (GWh/mo) & Post (GWh/mo) & $\Delta$ (GWh/mo) & Ratio \\
\midrule
CCGT & GE & 213 & 413 & +200 & 1.94 \\
 & GN & 397 & 932 & +534 & 2.34 \\
 & HC & 42 & 94 & +52 & 2.25 \\
 & IB & 146 & 316 & +169 & 2.16 \\
\addlinespace
Hydro & GE & -58 & -253 & -195 & 4.33 \\
 & GN & 5 & -33 & -38 & -6.96 \\
 & HC & -0 & -11 & -11 & 40.05 \\
 & IB & 33 & -84 & -117 & -2.55 \\
\addlinespace
Hydro\_pump & GE & 2 & -15 & -17 & -7.14 \\
 & GN & 4 & 3 & -1 & 0.70 \\
 & IB & 23 & -1 & -24 & -0.03 \\
\addlinespace
Nuclear & GE & 241 & 23 & -218 & 0.10 \\
 & IB & 1 & -137 & -138 & -- \\
\addlinespace
Wind & GN & -9 & -36 & -27 & 4.02 \\
 & HC & -13 & -50 & -36 & 3.71 \\
 & IB & 8 & -81 & -89 & -10.63 \\
\addlinespace
Solar PV & GE & -0 & -0 & -0 & 9.13 \\
 & GN & -3 & -11 & -8 & 3.29 \\
 & HC & 0 & 1 & +1 & 6.26 \\
 & IB & 3 & -42 & -45 & -12.10 \\
\addlinespace
\bottomrule
\end{tabular}
\end{table}

\textit{Reading.} Every CCGT firm doubles its post-DA gap (GN $534$~GWh/mo, the largest); Wind / Solar / Hydro / Nuclear mirror-decrease (REE substitutes CCGT for clean tech). Among Figure~\ref{fig:ree_intervention_full}'s codes, only code~61 (PO~3.2 system-security TR) is \emph{both} pay-as-bid \emph{and} local-grid binding --- the combination that yields zonal market power. Critical hours carry $3$--$5\times$ the per-day gap of flat or midday.

\subsection{Per-firm CCGT post-clearing revenue (proxy)}\label{sec:prov:phf_minus_pdbf:revenue}

\noindent Each (firm, tech, month) intervention multiplied by the regime-mean Fase~I-up pay-as-bid price (ESIOS id 705) gives a post-clearing redispatch revenue proxy (aggregates Fase~I+II + TR up).

\begin{table}[H]
\centering\scriptsize
\caption{\textbf{Per-firm CCGT up-redispatch revenue (proxy), EUR million per regime window.} Big-4 sum in DA15/ID15 ($\sim\texteuro 1{,}390$M) is close to the REE-published Fase~I-up cost ($\texteuro 1{,}200$M in Table~\ref{tab:cascade_costs_official}); the $\sim\!15\%$ overshoot reflects TR-up and Fase~II-up inclusion.}
\label{tab:ccgt_post_clearing_rev}
\setlength{\tabcolsep}{4pt}
% auto-built by scripts/analysis/regulatory/per_firm_RT2_revenue
% UP-redispatch revenue proxy = intervention_mwh × regime-mean Fase I up price (id 705)
% Captures Fase I + Fase II + TR up redispatch combined, not Fase I alone.
\begin{tabular}{l r r r r r r}
\toprule
Firm & 3-sess & ISP15-win & pre-blk & post-blk & \textbf{DA15/ID15} & Total \\
\midrule
\textbf{GN} & 336 & 337 & 84 & 751 & 703 & 2211 \\
\textbf{GE} & 195 & 164 & 32 & 331 & 343 & 1065 \\
\textbf{IB} & 145 & 84 & 27 & 230 & 286 & 771 \\
\textbf{HC} & 38 & 53 & 8 & 85 & 58 & 242 \\
\midrule
Big-4 sum & 714 & 638 & 151 & 1397 & 1390 & 4289 \\
\bottomrule
\end{tabular}

\end{table}

\textit{Reading.} \textbf{GN captures $\sim\!50\%$ of Big-4 CCGT post-clearing revenue in every regime.} GN CCGT: \texteuro 336M (3-sess) $\to$ \texteuro 337M (ISP15-win) $\to$ \texteuro 751M (post-blk) $\to$ \texteuro 703M (DA15/ID15) --- $\sim\texteuro 2$M/day pre-reforzada, $\sim\texteuro 5$--$8$M/day post --- concentrated in GN's southern-corridor zones (Sur/Levante/Cataluña; Table~\ref{tab:zonal_concentration}).

\subsection{CCGT marginal cost and the Fase~I scarcity rent}\label{sec:prov:phf_minus_pdbf:mc}

\noindent We compute a bottom-up CCGT marginal cost as $\mathrm{MC} = P_{\text{gas}}/\eta + P_{\text{CO}_2}\cdot\mathrm{EF}/\eta + \mathrm{VOM}$ with $\eta = 0.55$ LHV, $\mathrm{EF} = 0.20$~tCO$_2$/MWh-thermal, $\mathrm{VOM} = 4$~EUR/MWh-elec; gas from ESIOS id 1940, CO$_2$ from id 1391. Start-up/no-load adjustments shave $\sim\!10$~EUR/MWh off the premium below.

\begin{table}[H]
\centering\scriptsize
\caption{\textbf{Approximate CCGT marginal cost vs market and Fase~I prices, per regime} (EUR/MWh-elec). Premium $=$ Fase~I~up~$-$~MC$_{\text{central}}$ before the $\sim\!10$~EUR/MWh start-up adjustment.}
\label{tab:mc_ccgt_premium}
\setlength{\tabcolsep}{5pt}
\begin{tabular}{l r r r r}
\toprule
Regime & DA spot & Fase~I up (id 705) & MC$_{\text{central}}$ & Premium \\
\midrule
3-sess (Jun--Nov 24) & 77 & 159 & $\sim$100--115 & $44$ to $59$ \\
ISP15-win (Dec24--Mar25) & 118 & 187 & $\sim$125--150 & $37$ to $62$ \\
DA60/ID15 pre-blk & 54 & 145 & $\sim$130 & $15$ \\
DA60/ID15 post-blk & 63 & 165 & $\sim$100--110 & $55$ to $65$ \\
\textbf{DA15/ID15} & 76 & 168 & $\sim$100--110 & \textbf{$58$ to $68$} \\
\bottomrule
\end{tabular}
\end{table}

\textit{Reading.} Fase~I-up exceeds bottom-up CCGT MC by $40$--$70$~EUR/MWh in every regime: a \textbf{30--60\% mark-up over short-run MC} (after start-up adjustment) captured pay-as-bid by zonally-dominant CCGTs. DA spot in 2024--26 sits at or below CCGT MC much of the time. Scope: the units that clear Fase~I, disproportionately the zonally-constrained corridor CCGTs ($\sim\!20$--$30\%$ of the fleet, Table~\ref{tab:zonal_concentration}).


% =====================================================================
\section{Geographic CCGT concentration and zonal market power}\label{sec:prov:geo}
% =====================================================================

\textbf{Why this matters.} PO 3.2 redispatches relieve transmission-grid constraints \emph{locally}, and the resulting redispatch energy is paid pay-as-bid (\S\ref{sec:prov:phf_minus_pdbf}). The zonal-dominant CCGT operator therefore holds local market power on the Fase~I / TR bids in its zone, independent of the EUPHEMIA day-ahead clearing. The fewer the eligible CCGTs in a constrained zone, the higher the resulting Fase~I price.

\textbf{Variable.} CCGT installed capacity (MW) per (zone, firm), with HHI computed on capacity shares within each zone ($\times 10{,}000$). \textbf{Sources}: OMIE unit roster (zone $=$ plant city, conservative proxy for the REE PO 3.2 zonal book) joined with ESIOS archive 110 (master data on registered installations).

\begin{figure}[H]
\centering
\includegraphics[width=0.75\textwidth]{../../figures/working/fig_ccgt_zonal_capacity.pdf}
\caption{\textbf{Spanish CCGT installed capacity (MW) by zone and firm.} Bars stack each firm's MW; zone labels (Sur, Levante, etc.) are grouped from plant-city assignments. Galicia and Centro are duopolies; Sur / Levante / Cataluña have wider firm diversity but GN holds the largest share in each.}
\end{figure}

\begin{table}[H]
\centering
\caption{\textbf{Per-zone CCGT capacity concentration.} HHI on within-zone capacity shares ($\times 10{,}000$).}
\label{tab:zonal_concentration}
\begin{tabular}{l r r r l r r}
\toprule
Zone & Plants & Capacity (MW) & \#~firms & Top firm & Top share (\%) & HHI \\
\midrule
Galicia & 2 & 1\,247 & 2 & \textbf{GE} & 68.6 & 5\,695 \\
Norte & 10 & 2\,347 & 2$^\ast$ & non-Big-4$^\ast$ & 67.0 & 5\,577 \\
Aragón & 3 & 1\,870 & 2$^\ast$ & non-Big-4$^\ast$ & 57.7 & 5\,119 \\
Centro & 2 & 759 & 2 & \textbf{IB} & 50.9 & 5\,001 \\
Levante & 12 & 6\,117 & 3 & \textbf{GN} & 40.6 & 3\,616 \\
Cataluña & 8 & 3\,788 & 4 & \textbf{GN} & 44.5 & 3\,348 \\
Sur & 13 & 5\,563 & 4 & \textbf{GN} & 42.5 & 3\,040 \\
\bottomrule
\multicolumn{7}{l}{\footnotesize{$^\ast$~Norte and Aragón Big-4 share is small; top firms are non-Big-4 (Moeve, Engie, Total, Axpo).}} \\
\end{tabular}
\end{table}

\textit{Reading.} Galicia (HHI~5{,}695, 2 plants, GE~$69\%$) and Centro (HHI~5{,}001, 2 plants, IB~$51\%$/GN~$49\%$) are CCGT duopolies. The three southern-corridor zones --- Cataluña, Levante, Sur (HHI~$3{,}000$--$3{,}600$) --- have GN at $40$--$45\%$ and bear the heaviest reforzada post-blackout. Norte and Aragón are non-Big-4 dominated (Moeve, Engie, Total, Axpo).

\subsection{Bid-side heterogeneity by zonal dominance}\label{sec:prov:geo:bid_by_dominance}

\noindent We split each CCGT unit into ``dominant in its zone'' (largest within-zone capacity share) vs ``non-dominant'' and compare the median DA sell-bid price by firm.

\begin{table}[H]
\centering
\caption{\textbf{Median DA sell-bid price by firm and zonal dominance, EUR/MWh.} Per (unit-day) median; Pre $=$ 2024 same-calendar, Post $=$ 2025-05+ reforzada. Unit-counts (dominant / non-dominant): GN 15/2, GE 1/4, IB 1/9, HC 0/2.}
\label{tab:bid_by_zonal_dominance}
\setlength{\tabcolsep}{6pt}
\begin{tabular}{l r r r r r r r r}
\toprule
 & \multicolumn{4}{c}{Pre (2024 same-cal)} & \multicolumn{4}{c}{Post (2025-05+)} \\
 \cmidrule(lr){2-5} \cmidrule(lr){6-9}
Firm & Dom & Non-dom & Diff & ratio & Dom & Non-dom & Diff & ratio \\
\midrule
GN & 844 & 866 & $-22$ & 0.97 & 840 & 852 & $-12$ & 0.99 \\
GE & \textbf{2{,}680} & 2{,}381 & $299$ & 1.13 & \textbf{1{,}789} & 445 & \textbf{$1{,}344$} & \textbf{4.02} \\
IB & 156 & 170 & $-14$ & 0.92 & 147 & 196 & $-49$ & 0.75 \\
HC & -- & 156 & -- & -- & -- & 163 & -- & -- \\
\bottomrule
\end{tabular}
\end{table}

\textit{Reading.} \textbf{GN}: cap-padded $\sim\!\texteuro 840$/MWh ladder applied \emph{uniformly} across zones and regimes --- fleet-wide Cournot-on-quantity policy. \textbf{GE}: zone-specific extensive-margin withholding --- PGR5 in Galicia is fully parked, BES3/BES5/COL4 collapsed from parked to competitive post-reform. \textbf{IB} and \textbf{HC}: bid \texteuro 150--220/MWh near short-run marginal cost (\S\ref{sec:prov:phf_minus_pdbf:mc}); IB shows no dominant-vs-non-dominant differential. The CNMC's 2019 Naturgy sanction (\S\ref{sec:prov:cournot_cnmc:sanctions}) covered exactly the GN southern-corridor plants whose flat \texteuro 840 ladder appears above.

\subsection{Spatial test: did MTU15 swings predict the post-blackout restriction burden?}\label{sec:prov:geo:spatial_blackout}

\noindent CNMC's blackout report flags MTU15 as causing ``transitions between negative and positive prices [\dots] leading to strong changes in production volumes in certain zones''. We test bigger MTU15 within-hour swings $\to$ bigger post-blackout restrictions two ways. (a) \textbf{Per-zone BSTS for CCGT} on daily CCGT PHF $-$ PDBF (cutover 2025-04-28, Spec~A wind+solar+gas covariates + weekly cycle internally absorb seasonality). (b) \textbf{Per-unit cross-section} on all dispatchables (CCGT/Hydro/Pump), regressing the post-minus-pre PHF $-$ PDBF jump on the unit's within-hour SD with tech FE; \emph{seasonality is differenced out by same-calendar matching} (pre $=$ May--Dec 2024, post $=$ May--Dec 2025, same 8 months one year apart). Volatility score in both: per-unit within-hour SD of MW in the 40-day MTU15-IDA pre-blackout window (PHF).

\begin{table}[H]
\centering\scriptsize
\caption{\textbf{Spatial test results.} (a) Per-zone BSTS for CCGT: AbsEffect on daily PHF $-$ PDBF (GWh/day), pre 2024-06-14 to 2025-04-27, post 2025-04-28 to 2026-01-20. (b) Per-unit Spearman correlation of (post-pre) PHF $-$ PDBF jump with the unit's MTU15-IDA pre-blk within-hour MW SD, within tech.}
\label{tab:spatial_blackout}
\setlength{\tabcolsep}{4pt}
\begin{tabular}{l r r r @{\hspace{1cm}} l r r}
\toprule
\multicolumn{4}{c}{\textbf{(a) Per-zone CCGT BSTS}} & \multicolumn{3}{c}{\textbf{(b) Per-unit, within tech}} \\
\cmidrule(lr){1-4} \cmidrule(lr){5-7}
Zone & $n$ & vol (MW) & BSTS effect & Tech & $n$ & Spearman$_\rho$ \\
\midrule
Sur      & 9 & 4.09 & $\mathbf{15.29}$         & CCGT       & 33 & $-0.17$ \\
Levante  & 9 & 4.09 & $5.47$                   & Hydro      & 51 & $\mathbf{-0.57}$ \\
Cataluña & 5 & 1.70 & $\mathbf{4.97}$          & Hydro\_pump & 8 & $\mathbf{-0.86}$ \\
Norte    & 9 & 4.52 & $\mathbf{3.93}$          & \multicolumn{3}{l}{Pooled (tech FE):} \\
Aragón   & 2 & 18.75 & $\mathbf{3.37}$         & \multicolumn{3}{l}{$\beta=-0.046$ ($p\!<\!0.001$)} \\
Galicia  & 2 & 1.51 & $\mathbf{1.35}$          &       &      &      \\
Centro   & 2 & 3.65 & $0.75$                   &       &      &      \\
\bottomrule
\multicolumn{7}{l}{Cross-zone correlation (a): Pearson $-0.08$, Spearman $0.23$. Bold $\equiv$ 95\% CI excludes zero / $p\!<\!0.01$.}
\end{tabular}
\end{table}

\textit{Reading.} (i) The restriction jump is strongly heterogeneous across zones (Sur $15.3$ vs Centro $0.75$ GWh/day) --- the zonal dimension of reforzada is real --- and tracks the zonal-dominance map (\S\ref{sec:prov:geo}): the four highest-effect zones sit in the southern-Mediterranean voltage-fragile corridor. (ii) But \emph{volatility does not predict restriction burden in any tech} (after same-calendar seasonality control): the cross-zone CCGT correlation is weakly positive ($0.23$); the within-tech per-unit correlations for hydro and pump are sharply \emph{negative} ($-0.57$, $-0.86$); the pooled tech-FE slope is $-0.046$ ($p\!<\!0.001$). The CNMC's zone-specificity is empirically supported; their proposed channel (MTU15 swings drive restrictions) is falsified --- if anything, REE restricts the units that were \emph{steadier} under MTU15-IDA, suggesting the binding constraint is grid topology / voltage criticality, not within-hour swing magnitude. Reproduction in \texttt{scripts/analysis/regulatory/}.


% =====================================================================
\section{System cost cascade: the full ajuste view}\label{sec:prov:full_cascade}
% =====================================================================

\noindent Every MWh cleared in DA/IDA triggers downstream settlements in $\ge 6$ ancillary processes; the cascade is the falsifier for whether MTU15 reduces total system cost. The two pay-as-bid legs (Fase~I and TR) are where local market power can bite; the marginal-priced balancing services (aFRR PICASSO, mFRR MARI, RR TERRE) do not.

\subsection{Per-regime cost composition (REE-published cost indicators)}\label{sec:prov:full_cascade:costs}

\noindent Total cost (EUR million) per regime, summed across ISP cells. \textbf{Sources}: REE-published cost indicators 1373/1374 (Fase~I up/dn), 1375/1376 (Fase~II up/dn), 1723/1724 (TR up/dn), 712/2127 (aFRR reserve up/dn), 726/727 (imbalance excess/deficit).

\begin{table}[H]
\centering\scriptsize
\caption{\textbf{Per-regime ajuste-cascade cost composition.} EUR \emph{million} per regime, summed across ISP-15 cells. Negative TR (dn) values reflect REE paying providers to come down. Per-day evolution of the same components in Figure~\ref{fig:cascade_costs_timeseries}.}
\label{tab:cascade_costs_official}
\setlength{\tabcolsep}{4pt}
% auto-built by scripts/analysis/regulatory/ree_full_cascade_v2.py
\begin{tabular}{lrrrrrrrr}
\toprule
cat & Fase I (dn) & Fase I (up) & Fase II & Imbalance & TR (dn) & TR (up) & aFRR reserve & Total \\
regime &  &  &  &  &  &  &  &  \\
\midrule
Pre-IDA (Jan-Jun 24) & 10 & 822 & 141 & 0 & 0 & 501 & 208 & 1682 \\
3-sess (Jun-Nov 24) & 70 & 843 & 313 & 0 & -4 & 668 & 273 & 2163 \\
ISP15-win (Dec24-Mar25) & 33 & 573 & 239 & 0 & -1 & 627 & 264 & 1735 \\
DA60/ID15 pre-blk & 1 & 305 & 30 & 0 & -17 & 148 & 53 & 520 \\
DA60/ID15 post-blk & 110 & 1625 & 425 & 0 & -66 & 452 & 195 & 2741 \\
DA15/ID15 (Oct-Dec 25) & 44 & 1200 & 431 & 0 & -14 & 170 & 113 & 1944 \\
\bottomrule
\end{tabular}

\end{table}

\begin{figure}[H]
\centering
\includegraphics[width=0.78\textwidth]{../../figures/working/cascade_costs_timeseries.pdf}
\caption{\textbf{Daily REE ajuste-cost composition by service ($14$-day rolling mean), stacked positive.} All components are gross costs REE pays providers; UP and DN flows are stacked together with no sign convention. Vertical dotted lines mark the five reform dates (IDA reform 2024-06-14, ISP15 2024-12-11, MTU15-IDA 2025-03-19, blackout 2025-04-28, MTU15-DA 2025-10-01). The TR (up) channel falls sharply after MTU15-IDA, the granularity efficiency channel; Fase~I (up) balloons after the blackout, the reforzada-driven CCGT recall. Per-regime aggregates are in Table~\ref{tab:cascade_costs_official}.}
\label{fig:cascade_costs_timeseries}
\end{figure}

\subsection{Per-day equivalents (the headline)}\label{sec:prov:full_cascade:perday}
\begin{center}\scriptsize
\begin{tabular}{l r r r r r}
\toprule
 & 3-sess & ISP15-win & DA60/ID15 pre-blk & DA60/ID15 post-blk & \textbf{DA15/ID15} \\
Length (days) & $\sim$170 & $\sim$108 & 40 & $\sim$155 & 92 \\
Total cost (\texteuro M) & 2\,163 & 1\,735 & 520 & 2\,741 & 1\,944 \\
\textbf{Per day (\texteuro M)} & 12.7 & \textbf{16.1} & 13.0 & \textbf{17.7} & \textbf{21.1} \\
\textbf{Fase~I up / day} & 5.0 & 5.3 & 7.6 & 10.5 & \textbf{13.0} \\
\textbf{TR up / day} & 3.9 & 5.8 & 3.7 & 2.9 & 1.8 \\
\bottomrule
\multicolumn{6}{p{14cm}}{\scriptsize Raw per-day rates from Table~\ref{tab:cascade_costs_official}. Regime windows differ in calendar mix, so the row-by-row deltas mix granularity, reforzada, and seasonality; for a per-day evolution that lets the eye separate the reform-driven inflection from the seasonal noise, see Figure~\ref{fig:cascade_costs_timeseries}.} \\
\end{tabular}
\end{center}

\textit{Reading.} Total ajuste cost \texteuro 16.1M $\to$ \texteuro 21.1M/day ($31\%$) from ISP15-win to DA15/ID15: cost \emph{redistributed} from TR ($-\texteuro 4.0$M/day; TR price halved \texteuro 545 $\to$ \texteuro 226/MWh as MTU15-IDA absorbed intra-hour scarcity) to Fase~I ($\texteuro 5.7$M/day; price barely moved \texteuro 187 $\to$ \texteuro 168 but cost more than doubled). MTU15 is \emph{not} purely cost-saving; the $\sim\!\texteuro 21$M/day ajuste bill dwarfs any DA-auction margin. The TR price halving aligns with the IDA dual-direction surge documented in \S\ref{sec:prov:arbitrage_bsts}: finer-grained intraday bidding eroded the intra-hour scarcity that previously priced TR.

\subsection{Robustness panels}\label{sec:prov:full_cascade:robustness}

\noindent \emph{Balancing stack.} aFRR offer curve (ESIOS \texttt{curvas\_ofertas\_afrr}) and mFRR aggregated bids (ENTSO-E) show marginal-priced balancing is not concentration-distorted by the reforms ($60$--$90$ price steps per ISP, $\sim\!15\%$ MW growth, aFRR-Down price-range expansion \texteuro 115 $\to$ 384 tracking the Fase~I-dn growth in Table~\ref{tab:cascade_costs_official}; \texttt{scripts/analysis/balancing/balancing\_stack\_per\_regime.py}).
\emph{Cross-border price gap.} France is not a clean control for ID15 or DA15 (EU-wide SIDC/SDAC) but IS for Spain-specific ISP15 / reforzada. The Spain$-$France hourly DA-price gap widens against Spain by $\sim\!17$~EUR/MWh across the 40-day MTU15-IDA pre-blackout window --- ID15 is Spain-specific in magnitude despite the continental SIDC reform --- and shifts $37$~EUR/MWh at the reforzada onset (the Fase~I rent at the price level). Reproduction: \texttt{scripts/analysis/regulatory/cross\_border\_da\_price\_gap.py}; quantitative detail in \emph{preliminary results} \S 4.A.1 footnote.

% =====================================================================
\section{Cournot reading and CNMC enforcement precedent}\label{sec:prov:cournot_cnmc}
% =====================================================================

\begin{table}[H]
\centering\scriptsize
\caption{\textbf{CNMC sanction resolutions relevant to the practice.} SNC-DE-175-17 covers Besós 4, PBCN 1/2, Sagunto 1/2/3, Málaga 1, SROQ 1 --- the Cataluña / Levante / Sur zonal map (\S\ref{sec:prov:geo}); CNMC benefit estimate \texteuro 13.0~M over 4 months. LSE $=$ Ley del Sector Eléctrico 24/2013.}
\label{tab:cnmc_sanctions}
\setlength{\tabcolsep}{3pt}
\begin{tabular}{@{}l l l l r@{}}
\toprule
Year & Firm / Plant & Expediente & LSE article & Sanction \\
\midrule
2015 & Iberdrola hydro & SNC-DE-046-14 & 64 muy grave (manipulation) & 25.0 M EUR \\
2019 & \textbf{Naturgy CCGT} & SNC-DE-175-17 & 65.34 grave (abnormal DA offers) & $\sim$7 M EUR \\
2019 & Endesa Besós 3 & SNC-DE-174-17 & 65.34 grave (abnormal DA offers) & 6.6 M EUR \\
2023 & Naturgy Sabón 3 (Galicia) & --- & 65.33 grave (RT abuse) & confidential \\
2023 & Engie Castelnou & SNC-DE-152-22 & grave (availability) & --- \\
2023 & Ignis Escatrón 2 & SNC-DE-151-22 & grave (availability) & --- \\
\bottomrule
\end{tabular}
\end{table}

\noindent GN's near-constant $\sim\!\texteuro 840$/MWh DA ladder (Table~\ref{tab:bid_by_zonal_dominance}) with quantity-side adjustment is the signature of Cournot-on-quantity --- price-raising in DA dressed as ``my offer price has not changed,'' with rent captured in pay-as-bid Fase~I (\S\ref{sec:prov:phf_minus_pdbf:revenue}). The CCGT $\sigma_p$ widening of \emph{preliminary results} \S A.4 sits on the in-band competing tier, not on this parked ladder. The 2019 Naturgy sanction (SNC-DE-175-17) is the precedent; 2024--26 is the same at fleet scale, with granularity as an extra tool.\footnote{The $\sim\!56$ April-2026 voltage-control expedientes (Art~65.8) are a separate regulatory channel.}

\section{Do unit-level arbitrageurs exist? Buy and sell offers from production plants}\label{sec:prov:arbitrage_units}

The theoretical model in \emph{model.tex} posits an arbitrageur agent that links the day-ahead (DA) and intraday auction (IDA) prices via a position $s$. Empirically, the natural signature of unit-level arbitrage is that a \emph{production} unit (one that generates electricity) submits not only sell offers (the default for a producer) but also buy offers, i.e.\ repurchases part of an earlier sell-side commitment. We test this directly in OMIE offer-header data (\texttt{cab} for DA, \texttt{icab} for IDA, 2018-01-01 to 2026-01-31). For each unit-day (DA) or unit-session (IDA), we record whether at least one sell offer (\texttt{buy\_sell = 'V'}) and at least one buy offer (\texttt{buy\_sell = 'C'}) appear; we then aggregate by tech-group and report the share of seller unit-days/sessions that also carry a buy offer (Table~\ref{tab:arbitrage_units}).

\begin{table}[h]
\centering
\footnotesize
\caption{Production-unit buy-and-sell offers in OMIE auctions, by tech-group, 2018-01-01 to 2026-01-31. ``\% with both'' is the share of unit-days (DA) or unit-sessions (IDA) on which the unit submitted at least one sell offer that \emph{also} carried at least one buy offer. The DA market is pure sell-side for production: no unit ever submits a DA buy offer. IDA flips: every tech-group has non-trivial dual-direction bidding, with variable RES the most extreme.}
\label{tab:arbitrage_units}
\begin{tabular}{l r r r r r}
\toprule
Tech-group & n unit-days (DA) & \% with both (DA) & n unit-sessions (IDA) & \% with both (IDA) \\
\midrule
Solar PV   &   944{,}918 & 0.00 & 2{,}214{,}353 & 30.39 \\
Wind       &   537{,}764 & 0.00 & 2{,}389{,}721 & 17.62 \\
Biomass    &   149{,}641 & 0.00 &   116{,}343 &  7.23 \\
Hydro      &   119{,}357 & 0.00 &   347{,}350 &  6.37 \\
Hydro\_pump &   21{,}533 & 0.00 &    54{,}949 &  5.09 \\
Nuclear    &    42{,}769 & 0.00 &    13{,}799 &  4.44 \\
Hydro\_RES  &  311{,}347 & 0.00 &   446{,}228 &  4.38 \\
Cogen      &   511{,}140 & 0.00 &   345{,}407 &  4.33 \\
CCGT       &   137{,}481 & 0.00 &   271{,}139 &  4.25 \\
Coal       &     5{,}320 & 0.00 &     7{,}119 &  0.81 \\
\bottomrule
\end{tabular}
\end{table}

\noindent Two facts stand out. First, the DA market is purely sell-side for production: across more than 2.7 million unit-days and ten tech-groups we observe \emph{zero} DA buy offers from any production unit. This is consistent with DA being settled as a single hourly product per unit per day --- there is no within-day price spread for a producer to exploit by buying back its own forward position, and the EUPHEMIA-style auction does not naturally accommodate units submitting both directions. Second, the IDA picture is the opposite: every tech-group has at least some unit-sessions with dual-direction bidding. Variable renewable sources (Solar PV at 30\%, Wind at 18\%) dominate --- consistent with forecast-error rebalancing, where the operator buys back an earlier sell when the realised production forecast tightens. Among dispatchable techs, CCGT (4.3\%), Hydro\_pump (5.1\%), Nuclear (4.4\%) and Hydro (6.4\%) all show non-zero shares; only Coal, with a near-zero share (0.81\%), behaves as a pure seller.

\noindent The reading consistent with the model is the following. The DA's auction structure (single price per unit per hour, EUPHEMIA-style) leaves no within-day spread for a producer to arbitrage in DA itself; arbitrage between DA and IDA is therefore observed not in the DA bid file but in the IDA bid file, where the producer rebalances. Variable-RES techs do most of this rebalancing because their forecast updates are largest; thermal and Hydro\_pump do less but the share is not zero, which is direct evidence that the arbitrageur agent in the model is empirically active. Moreover, the share is bounded in the IDA --- at most $\sim$30\% even for Solar PV --- consistent with capacity-binding limited arbitrage (Result~3 of Ito and Reguant, 2016), not the unlimited-strategic regime in which every product would carry an arbitrage position.

\subsection{Measure construction: what ``arbitrage intensity'' captures and what it doesn't}\label{sec:prov:arbitrage_measure}

The measure used in this section requires careful documentation because it is not a direct observation of arbitrage profit, capital, or net position. We construct it as follows:

\textbf{Definition.} For each calendar date $d$, IDA session $s$, and production unit $u$, we record whether the unit submitted at least one sell offer (\texttt{buy\_sell = 'V'}) \emph{and} at least one buy offer (\texttt{buy\_sell = 'C'}) in the IDA offer-header file \texttt{icab\_all}. A unit-session is a \emph{dual-direction} bid if both directions appear. Our daily arbitrage-intensity outcome is the per-tech-group count of dual-direction unit-sessions on date $d$, summed across the three IDA sessions of the day.

\textbf{What the measure captures.} A unit-session with both buy and sell offers is consistent with three economically distinct behaviours, all of which contain an arbitrage component in the broad sense: (i)~\emph{position rebalancing}, where the unit's DA-side commitment exceeds (or falls short of) its updated production forecast and the unit submits a buy (sell) to close the gap; (ii)~\emph{pure financial arbitrage}, where the unit buys at a low spot price expecting to sell at a higher spot price (or to net out against its DA position) without any underlying production change; and (iii)~\emph{strategic withholding then re-clearing}, where the unit deliberately under-offers at DA expecting to capture a higher IDA clearing price, then re-buys some of the previously-sold MW. All three are forms of unit-level arbitrage between the DA-cleared price and the IDA-cleared price.

\textbf{What the measure does NOT capture.} The measure is binary at the unit-session level: it does not encode the \emph{volume} of dual-direction offers (a 1~MW buy and a 100~MW sell count the same as a 50/50 split), the \emph{price-spread} between the buy and sell offers, or the \emph{net} cleared MWh direction across the session. It also does not capture cross-session arbitrage (selling in IDA-1, buying in IDA-2), within-day arbitrage of the same unit's IDA position, or arbitrage between IDA and the continuous-intraday market XBID. Within-session netting between sell and buy offers can occur in our measure even when total unit MWh stays unchanged, so high dual-direction counts can co-exist with low absolute net repositioning. Conversely, a unit shifting its entire 100~MW sell from session~1 to session~2 (no buys) is genuine arbitrage that our measure misses entirely.

\textbf{Sources of bias for the model interpretation.} The model's arbitrageur agent is profit-motivated (carries position $s$ to exploit the forward-spot wedge). Our measure conflates this with two innocent behaviours: forecast updates (especially for RES, which see large within-hour forecast revisions) and operational re-positioning around physical constraints. The Solar PV 30\% share probably reflects forecast updates more than profit-driven arbitrage; the much lower CCGT/Nuclear shares ($\sim$4\%) are closer to the model's intended profit-arbitrage activity, since thermal units have small forecast errors and minimal physical re-positioning needs. The measure should therefore be read as a \emph{lower} bound on the share of unit-sessions with any arbitrage component, not as a measure of pure profit-arbitrage intensity.

\textbf{Implication for the model's P7 prediction.} The cross-tech heterogeneity discussed below should be read against this background: the measure is well-suited to detect changes in dual-direction prevalence, but the level interpretation requires care.

\subsection{Did arbitrage intensity rise at MTU15-IDA? BSTS on daily dual-direction counts per tech}\label{sec:prov:arbitrage_bsts}

If the model's arbitrageur agent is empirically active, MTU15-IDA (cutover 2025-03-19) should mechanically raise the per-tech daily dual-direction count: each clock-hour now hosts four matched 15-min IDA products against one DA product, so the per-clock-hour space for arbitrage is four times larger. We run a per-tech BSTS on the daily dual-direction count, with wind+solar+gas covariates (matching \emph{preliminary results} Spec~A), pre-window 2022-06-14 (post 6$\to$3 IDA reform) to 2025-03-18, and the 40-day pre-blackout post-window 2025-03-19 to 2025-04-27. Cleanly identifies the MTU15-IDA effect with reforzada held constant (reforzada onset is 2025-04-28).

\begin{table}[h]
\centering
\footnotesize
\caption{BSTS effect of MTU15-IDA on daily dual-direction unit-session count, by tech-group. Stars on $p$-values: $^{*}p{<}0.10$, $^{**}p{<}0.05$, $^{***}p{<}0.01$. Pre mean and post mean are daily counts averaged over their respective windows; absolute effect is the BSTS counterfactual gap. The heterogeneity across techs is the headline read: forecast-driven techs (Solar PV, Hydro\_RES, Biomass) jump sharply, while thermal and pumping techs do not.}
\label{tab:bsts-arbitrage}
\begin{tabular}{l r r r r r}
\toprule
Tech-group & Pre mean & Post mean & Abs.\ effect & Rel.\ effect & $p$ \\
\midrule
Solar PV        &  828.8 & 1100.2 & $+271.4^{***}$ & $+33.6\%$ & $<\!0.001$ \\
Hydro\_RES      &   25.3 &   43.5 &  $+18.1^{***}$ & $+76.3\%$ & $<\!0.001$ \\
Biomass         &    3.2 &    9.8 &   $+6.6^{***}$ & $+521\%$  & $<\!0.001$ \\
Nuclear         &    0.4 &    2.0 &   $+1.6^{***}$ & $+450\%$  & $<\!0.001$ \\
\addlinespace
Wind            &  648.7 &  682.2 &  $+33.5$       & $+6.4\%$  & $0.32$ \\
Hydro           &   33.8 &   33.6 &   $-0.1$       & $+1.5\%$  & $0.48$ \\
Cogen           &   19.4 &   21.5 &   $+2.1$       & $+13.4\%$ & $0.22$ \\
\addlinespace
CCGT            &    6.1 &    5.1 &   $-1.0$       & $-27.1\%$ & $0.44$ \\
Hydro\_pump     &    5.6 &    4.4 &   $-1.2^{*}$   & $-19.7\%$ & $0.05$ \\
Coal            &    0.5 &    0.1 &   $-0.4^{***}$ & $-73.5\%$ & $0.001$ \\
\bottomrule
\end{tabular}
\end{table}

\noindent The pattern is sharply stratified by tech. \emph{Forecast-driven techs} (Solar PV, Hydro\_RES, Biomass) show massive jumps in dual-direction count, with Solar PV alone adding $\sim$270 dual-direction unit-sessions per day --- consistent with the mechanism in §4.A.3 of the preliminary results: per-period IDA scheduling lets aggregators track their 15-min production forecast more closely, which mechanically requires more buy-then-sell rebalancing as the forecast updates within the hour. \emph{Wind, Hydro and Cogen} are flat, plausibly because their dual-direction shares were already saturated before MTU15-IDA. \emph{Thermal and storage} (CCGT, Coal, Hydro\_pump) show no increase or a small decrease.

\noindent The CCGT and Hydro\_pump results require an extra layer of caveat before any model interpretation. By the MTU15-IDA post-window (March--April 2025), Spanish solar was crowding the merit order and CCGT day-ahead cleared volume had already collapsed to near zero (the $q_1$ drop in \S\ref{sec:prov:q1_drop}); a unit producing essentially nothing has very little to rebalance, so the CCGT BSTS null is largely a no-production result rather than a tested-and-rejected model prediction. The dual-direction count is a sensible \emph{rebalancing} measure only for techs that are actually clearing volume in the post-window --- which forecast-driven RES and biomass are (high cleared MW under MTU15-IDA), and which CCGT, Coal and Hydro\_pump are not. Pure profit-arbitrage by thermal units may also operate through net-direction changes across sessions (a phenomenon the measure misses, per \S\ref{sec:prov:arbitrage_measure}). The CCGT scale-up at DA15 documented in \emph{preliminary results} \S 4.A.2 --- which lands in a window where CCGT is producing again --- is the more direct empirical signature of thermal arbitrage, through quantity migration into DA rather than within-session position rebalancing in IDA. \emph{There is also a third arbitrage channel the within-OMIE dual-direction measure cannot see by construction:} during reforzada (continuous from 2025-04-28, overlapping all of the post-blackout window), the dominant CCGT arbitrage is between Fase~I up-redispatch (pay-as-bid, REE-paid) and the wholesale auction clear --- the parked-tier / Fase~I-recall pattern documented in \S\ref{sec:prov:phf_minus_pdbf}. \textbf{Read together, the three findings characterise three distinct empirical arbitrage channels}: (a)~\emph{within-session forecast-driven rebalancing in IDA}, dominated by RES/biomass and amplified by MTU15-IDA; (b)~\emph{between-market quantity migration into DA}, dominated by CCGT and amplified by DA15; (c)~\emph{Fase~I-vs-wholesale arbitrage under reforzada}, dominated by CCGT and driven by the post-blackout regulatory regime rather than by the granularity reforms.

\subsection{What to elevate from this document into preliminary\_results: selective glossary}\label{sec:prov:elevation_glossary}

The preliminary-results main memo carries the headline identification claims (Spec~A BSTS, Spec~B per-hour-class BSTS, Spec~C per-curve DiD). The findings in this companion document deepen those headlines but most do not warrant elevation. Below is a deliberately short list of items that I believe should move into the preliminary memo --- to its appendix unless flagged otherwise --- because they discipline a model claim or close an identification loop the preliminary memo currently leaves open. Items not on this list stay here.

\begin{itemize}\setlength{\itemsep}{4pt}
\item \textbf{Production-unit arbitrageurs exist (§\ref{sec:prov:arbitrage_units})} $\to$ preliminary memo \emph{appendix}. Directly validates the model's arbitrageur agent (P7): production units submit zero DA buy offers but 4--30\% dual-direction IDA unit-sessions across all tech groups. Worth a half-page appendix figure plus the per-tech table.
\item \textbf{Dual-direction count jumps at MTU15-IDA for forecast-driven techs (§\ref{sec:prov:arbitrage_bsts})} $\to$ preliminary memo \emph{appendix}. Direct identification of one arbitrage channel via BSTS. Pair with the §4.A.3 renewable-cleared-MW reallocation finding (which it complements, not duplicates).
\item \textbf{Geographic CCGT concentration and zonal market power (§\ref{sec:prov:geo})} $\to$ preliminary memo \emph{main text} (one paragraph) or \emph{appendix}. Provides the market-structure scaffolding for the §4.A.2 CCGT scale-up reading: the bidder is identifiably concentrated in specific zones, which makes the strategic-bidder identification of the model concrete in the data.
\item \textbf{Cournot reading and CNMC enforcement precedent (§\ref{sec:prov:cournot_cnmc})} $\to$ preliminary memo \emph{main text} (one sentence) or footnote. The 2019 Naturgy sanction (SNC-DE-175-17) is the regulatory analogue of the model's strategic-bidder withholding mechanism. Useful for the §4.A.2 ``competitive clearing anticipates expected REE's intervention'' footnote.
\end{itemize}

\noindent The list is intentionally short: four items in total, two of which (the arbitrage findings) come from this session's work and the other two of which (geographic concentration, CNMC precedent) discipline existing headline claims rather than introducing new ones. The bid-curve-shape decompositions (§\ref{sec:prov:dw_decomp}), tech-mix by hour-class (§\ref{sec:prov:tech_mix}), continuous-intraday repositioning (§\ref{sec:prov:ci_per_tech}), $q_1$ CCGT drop (§\ref{sec:prov:q1_drop}), REE post-clearing inclusion (§\ref{sec:prov:phf_minus_pdbf}), and full ajuste cascade (§\ref{sec:prov:full_cascade}) all serve as supplements to existing preliminary-memo findings rather than as candidates for elevation in their own right.

\end{document}
